\chapter[Przestrzenie Sobolewa]{Przestrzenie Sobolewa: słabe pochodne, aproksymacja i przedłużanie}
\label{chap:pde-sobolev}

Na prostej funkcja $u(x)=|x|$ nie jest różniczkowalna w zerze,
lecz jej pochodna w sensie dystrybucji jest funkcją
$\operatorname{sgn}x$. Z drugiej strony pochodna funkcji
Heaviside'a $H=\chi_{(0,\infty)}$ jest deltą Diraca,
a nie funkcją lokalnie całkowalną (sekcja
~\ref{sec:pde7-derivatives}). Ta różnica motywuje przestrzenie Sobolewa: dopuszczamy
niegładkie funkcje, ale wymagamy, aby ich pochodne dystrybucyjne
były reprezentowane przez funkcje z określonych klas całkowalności.

Niech w całym rozdziale $d\ge1$ będzie liczbą
całkowitą, $k\in\N_0$, $\Omega\subset\R^d$ niepustym zbiorem
otwartym, a $1\le p\le\infty$. Nie zakładamy spójności $\Omega$,
o ile nie wskażemy jej jawnie. Chyba że wyraźnie zaznaczono
inaczej, funkcje są rzeczywiste. Elementy przestrzeni $L^p$
są klasami funkcji równych prawie wszędzie. Pochodne rozumiemy
w sensie dystrybucji, a ich reprezentanty funkcyjne i równości
w $L^p$ — z dokładnością do zbiorów miary zero. Wszystkie całki na $\Omega$ są liczone
względem miary Lebesgue'a $\lambda_d$. Definicje dystrybucji i
pochodnych słabych znajdują się w rozdziale~\ref{chap:pde-distributions},
a przestrzenie $H^s(\R^d)$ w sekcji~\ref{sec:pde8-bessel}.

\section{Definicje i pierwszy słownik}
\label{sec:pde9-def}

\begin{definicja}[Przestrzenie $W^{k,p}$]
\label{def:pde9-Wkp}
Dla $k\in\N_0$ i $1\le p\le\infty$ definiujemy
\[
 W^{k,p}(\Omega)=\{u\in L^p(\Omega):
  \partial^\alpha u\in L^p(\Omega)
  \text{ dla każdego }\alpha\in\N_0^d,\ |\alpha|\le k\},
\]
gdzie $\partial^\alpha u$ oznacza pochodną dystrybucyjną.
Każda funkcja $u\in L^p(\Omega)$ jest lokalnie całkowalna:
dla zwartego $K\subset\Omega$ nierówność Höldera daje
$\int_K |u|\,\mathrm d\lambda_d
\le\lambda_d(K)^{1-1/p}\|u\|_p$ dla $1<p<\infty$;
dla $p=1$ wystarcza całkowalność globalna,
a dla $p=\infty$ oszacowanie
$\lambda_d(K)\|u\|_\infty$. Definiuje więc dystrybucję regularną.
Wyrażenie $\partial^\alpha u\in L^p$ oznacza istnienie
funkcji $v_\alpha\in L^p(\Omega)$ spełniającej
\[
 \int_\Omega u(x)\partial^\alpha\varphi(x)\,\mathrm d\lambda_d(x)
 =(-1)^{|\alpha|}\int_\Omega v_\alpha(x)\varphi(x)\,
 \mathrm d\lambda_d(x)
 \qquad(\varphi\in C_c^\infty(\Omega)).
\]
Zgodnie z definicją~\ref{def:pde7-weak} i jednoznacznością
reprezentantów dystrybucji regularnych (sekcja~\ref{sec:pde7-def})
$v_\alpha$ jest wyznaczona prawie wszędzie.
Dla $\alpha=0$ mamy $v_0=u$. Normę określamy wzorami
\begin{align*}
 \|u\|_{W^{k,p}(\Omega)}
  &=\left(\sum_{|\alpha|\le k}\|\partial^\alpha u\|_{L^p(\Omega)}^p
    \right)^{1/p},&&1\le p<\infty,\\
 \|u\|_{W^{k,\infty}(\Omega)}
  &=\max_{|\alpha|\le k}\|\partial^\alpha u\|_{L^\infty(\Omega)}.
\end{align*}
Dla $k=0$ otrzymujemy $W^{0,p}=L^p$, z równością norm.
Piszemy $H^k(\Omega)=W^{k,2}(\Omega)$,
a w szczególności $H^1(\Omega)=W^{1,2}(\Omega)$.
\end{definicja}

Liniowość pochodnych dystrybucyjnych oraz przestrzeni $L^p$
pokazuje, że $W^{k,p}(\Omega)$ jest przestrzenią wektorową.
Niech $M=\#\{\alpha\in\N_0^d:|\alpha|\le k\}$.
Odwzorowanie $u\mapsto(\partial^\alpha u)_{|\alpha|\le k}$
jest liniowe i iniektywne, gdyż zawiera współrzędną $u$.
Dla $p<\infty$ podana norma jest normą $\ell^p$
wektora norm składowych, a dla $p=\infty$ ich maksimum.
Nierówność trójkąta wynika z nierówności trójkąta w $L^p$
dla każdej składowej, a następnie z nierówności Minkowskiego
dla skończonych ciągów (dla $p=\infty$ z nierówności dla maksimum).
Jednorodność wynika z liniowości pochodnych, a zerowa norma
pociąga $\|u\|_p=0$. Zatem istotnie otrzymujemy normę.
Zupełność wykażemy w następnej sekcji.

\begin{definicja}[Przestrzeń lokalna i domknięcie funkcji testowych]
\label{def:pde9-local-zero}
Niech $k\in\N_0$, $1\le p\le\infty$.
Klasa funkcji mierzalnych $u$ należy do
$W^{k,p}_{\mathrm{loc}}(\Omega)$, jeśli
$u|_U\in W^{k,p}(U)$ dla każdego niepustego otwartego $U$
takiego, że $\overline U$ jest zwartym podzbiorem $\Omega$
(zapisujemy $U\Subset\Omega$).
Dla $1\le p<\infty$ definiujemy
\[
 W^{k,p}_0(\Omega)=\overline{C_c^\infty(\Omega)}^{\,W^{k,p}(\Omega)}.
\]
Indeks $0$ oznacza w tej chwili \emph{definicję przez domknięcie},
a nie jeszcze twierdzenie o istnieniu śladu na brzegu.
W tym rozdziale zapis $W^{k,p}_0$ stosujemy tylko dla $p<\infty$;
nie przenosimy twierdzeń o gęstości w tych normach na $p=\infty$.
\end{definicja}

Ograniczenie funkcji z $W^{k,p}(\Omega)$ do otwartego
$U\subset\Omega$ należy do $W^{k,p}(U)$ i
$\partial^\alpha(u|_U)=(\partial^\alpha u)|_U$:
test z $C_c^\infty(U)$ przedłużony przez zero jest testem
na $\Omega$, więc spełnia tę samą tożsamość całkową.
Norma po ograniczeniu nie rośnie, ponieważ całkujemy
na podzbiorze (lub bierzemy istotne supremum na podzbiorze).
W szczególności $W^{k,p}(\Omega)\subset W^{k,p}_{\mathrm{loc}}(\Omega)$.

W definicji lokalnej reprezentanty pochodnych z dwóch
takich otwartych zbiorów zgadzają się prawie wszędzie
na ich przecięciu, dzięki jednoznaczności dystrybucji regularnych.
Wybierzmy przeliczalne pokrycie $(B_n)_{n\ge1}$
zbioru $\Omega$ kulami o domknięciach zwartych w $\Omega$
(można użyć kul o wymiernych środkach i promieniach).
Niech $v_{\alpha,n}$ oznacza mierzalnego reprezentanta
pochodnej na $B_n$. Na mierzalnych, parami rozłącznych
zbiorach $A_n=B_n\setminus\bigcup_{j<n}B_j$
określamy $v_\alpha=v_{\alpha,n}$.
Zbiory $A_n$ pokrywają $\Omega$, więc $v_\alpha$ jest
mierzalna. Dla ustalonego $m$ na każdym $A_n\cap B_m$
zgadza się prawie wszędzie z $v_{\alpha,m}$, a
przeliczalna suma zbiorów wyjątkowych ma miarę zero.
Wobec tego $v_\alpha|_{B_m}=v_{\alpha,m}$ prawie wszędzie.
Każdy zwarty podzbiór $\Omega$ jest zawarty w skończonej
sumie tych kul, co zapewnia lokalną przynależność do $L^p$,
również dla $p=\infty$ przez maksimum skończenie wielu norm.
Dla dowolnego testu na $\Omega$ wybieramy otwarty zbiór
$U\Subset\Omega$ zawierający jego nośnik (na przykład
skończoną sumę tych kul). Reprezentant słabej pochodnej
na $U$ zgadza się prawie wszędzie z $v_\alpha$
na każdym $U\cap B_n$, więc również na całym $U$.
Tożsamość testowa na $U$ zachodzi z $v_\alpha$
i jest zarazem tożsamością na $\Omega$, ponieważ
test i jego pochodne znikają poza $U$. Zatem $v_\alpha$ reprezentuje
$\partial^\alpha u$ na całej $\Omega$.
Nie wymaga to globalnej całkowalności na $\Omega$.

\begin{uwaga}[Wybór normy]
Dla $M$ składowych i $1\le p<\infty$ nierówność
Höldera dla skończonych ciągów daje
\[
 \|u\|_{W^{k,p}}\le\sum_{|\alpha|\le k}\|\partial^\alpha u\|_p
 \le M^{1-1/p}\|u\|_{W^{k,p}}.
\]
Pierwsza nierówność wynika z tego, że dla nieujemnych
$a_i$, przy $S=\sum_i a_i>0$, mamy
$0\le a_i/S\le1$, więc $(a_i/S)^p\le a_i/S$.
Sumując i mnożąc przez $S^p$, otrzymujemy
$\sum_i a_i^p\le S^p$; dla $S=0$ jest to równość.
Dla $p=\infty$ analogicznie maksimum jest nie większe
od sumy, a suma nie przekracza $M$ razy maksimum.
Normy te są więc równoważne.
Przy $p=2$ norma $H^1(\Omega)$ spełnia
$\|u\|_{H^1}^2=\|u\|_2^2+\|\nabla u\|_2^2$.
Zgodnie z twierdzeniem~\ref{thm:pde8-Hs}, na całej $\R^d$
przestrzenie $H^k$ zdefiniowane przez słabe pochodne pokrywają się
z przestrzeniami Bessela z rozdziału~\ref{chap:pde-fourier-Rd};
normy są równoważne, a dla $k=1$ (oraz $k=0$)
są równe. Istotnie, dla $u\in W^{1,2}(\R^d)$ reguła
różniczkowania transformaty i twierdzenie Plancherela
z rozdziału~\ref{chap:pde-fourier-Rd} dają
\[
 \begin{aligned}
 \|u\|_{W^{1,2}}^2
 &=\|u\|_2^2+\sum_{j=1}^d\|\partial_j u\|_2^2\\
 &=(2\pi)^{-d}\int_{\R^d}(1+|\xi|^2)|\widehat u(\xi)|^2
 \,\mathrm d\lambda_d(\xi).
 \end{aligned}
\]
Prawa strona jest kwadratem normy Bessela $H^1$.
\end{uwaga}

\begin{przyklad}[Załamanie nie jest skokiem]
Niech $I=(-1,1)$. Funkcja $u(x)=|x|$ należy do $W^{1,p}(I)$
dla każdego $1\le p\le\infty$, ponieważ
$u'=\operatorname{sgn}x$ prawie wszędzie i pochodna ta jest
jednocześnie pochodną dystrybucyjną, jak wykazano w sekcji
~\ref{sec:pde7-derivatives}. Obie funkcje są ograniczone na $I$
o skończonej mierze, zatem należą do każdego z rozważanych
$L^p(I)$. Natomiast funkcja
$H=\chi_{(0,1)}$ nie należy do $W^{1,p}(I)$ dla żadnego
$1\le p\le\infty$: jej pochodna dystrybucyjna to $\delta_0$
(przykład~\ref{ex:pde7-heaviside}). Delta nie jest reprezentowana
nawet przez funkcję lokalnie całkowalną na $I$:
dowód z sekcji~\ref{sec:pde7-def}, wykorzystujący testy
o nośnikach kurczących się do zera, działa również na $I$.
W szczególności nie może być reprezentowana przez funkcję z $L^p(I)$.
Ponadto $u(x)=|x|$ nie należy do $W^{2,p}(I)$ dla żadnego
$1\le p\le\infty$, ponieważ $\partial^2 u=2\delta_0$
(co wynika z $\operatorname{sgn}x=2H(x)-1$ prawie wszędzie).
\end{przyklad}

\begin{przyklad}[Osobliwość potęgowa]
\label{ex:pde9-power}
Niech $d\ge2$, $0<a<d$, a $\psi\in C_c^\infty(\R^d)$
będzie równa $1$ w otoczeniu zera. Niech $u(x)=\psi(x)|x|^{-a}$
dla $x\ne0$ i $u(0)=0$ (wartość w zerze nie wpływa
na klasę funkcji). Dla $1\le p<\infty$ zachodzi
\[
 u\in L^p(\R^d)\iff ap<d,\qquad
 u\in W^{1,p}(\R^d)\iff (a+1)p<d.
\]
Dla $p=\infty$ funkcja $u$ nie należy nawet do $L^\infty$.

Istotnie, wybieramy $\delta>0$ tak, aby $\psi=1$ na
$B_\delta(0)$. Poza tą kulą funkcja i jej klasyczne pochodne
są gładkie, ograniczone i mają zwarty nośnik w rozważanym
obszarze, więc nie sprawiają problemu z całkowalnością.
Na kuli całkowanie radialne daje
\[
 \int_{B_\delta(0)} |u(x)|^p\,\mathrm d\lambda_d(x)
 =\omega_d\int_0^\delta r^{d-1-ap}\,\mathrm d\lambda_1(r),
\]
gdzie $\omega_d>0$ jest polem powierzchni sfery jednostkowej.
Istotnie, dla $0<\varepsilon<\delta$ całka
$\int_\varepsilon^\delta r^{d-1-ap}\,\mathrm d\lambda_1(r)$
jest równa $(\delta^{d-ap}-\varepsilon^{d-ap})/(d-ap)$
dla $ap\ne d$, a $\log(\delta/\varepsilon)$ dla $ap=d$.
Granica przy $\varepsilon\downarrow0$ jest skończona
dokładnie dla $ap<d$. Przy $ap=d$ rozbieżność jest
logarytmiczna, a przy $ap>d$ potęgowa. Ponadto na przekłutej kuli
$|u(x)|=|x|^{-a}$ przekracza dowolne ustalone dodatnie
ograniczenie na kuli o dodatniej mierze, co wyklucza $L^\infty$.

Dla $x\ne0$ klasyczny gradient jest równy
\[
 v(x)=|x|^{-a}\nabla\psi(x)-a\psi(x)|x|^{-a-2}x,
 \qquad |v(x)|=a|x|^{-a-1}\quad(0<|x|<\delta).
\]
Możemy przyjąć $v(0)=0$. Wszystkie składowe $v$ należą
do $L^p$ dokładnie wtedy, gdy $|v|\in L^p$:
korzystamy z $|v_j|\le |v|\le\sum_j|v_j|$ oraz
nierówności trójkąta w $L^p$. Na kuli
\[
 \int_{B_\delta(0)} |v(x)|^p\,\mathrm d\lambda_d(x)
 =a^p\omega_d\int_0^\delta r^{d-1-(a+1)p}\,
 \mathrm d\lambda_1(r).
\]
Stosując poprzedni rachunek z $a+1$ zamiast $a$,
otrzymujemy warunek skończoności $(a+1)p<d$;
na granicy ponownie występuje rozbieżność logarytmiczna.

Załóżmy $(a+1)p<d$. Wówczas $u,v_j\in L^p$ oraz
$a<d-1$ (ponieważ $p\ge1$). Dla dowolnego
$\varphi\in C_c^\infty(\R^d)$ wybieramy $R>\delta$
tak duże, aby jej nośnik był zawarty w $B_R(0)$.
Na pierścieniu $B_R(0)\setminus\overline{B_\varepsilon(0)}$,
$0<\varepsilon<\delta$, całkowanie przez części daje
\[
 \int_{B_R\setminus\overline{B_\varepsilon}}
 u\,\partial_j\varphi\,\mathrm d\lambda_d
 =-\int_{B_R\setminus\overline{B_\varepsilon}}
 v_j\varphi\,\mathrm d\lambda_d
 +\int_{\partial B_\varepsilon}
 u\varphi\,\nu_j\,\mathrm d\sigma,
\]
gdzie $\nu=-x/|x|$ jest normalną zewnętrzną pierścienia
na wewnętrznej sferze, a $\sigma$ jest miarą powierzchniową
tej sfery. Na zewnętrznej sferze test znika.
Moduł składnika brzegowego jest nie większy niż
$\|\varphi\|_\infty\omega_d\varepsilon^{d-1-a}$,
co dąży do zera. Obie całki objętościowe zbiegają do
całek na $\R^d$ z twierdzenia o zbieżności zmajoryzowanej:
majorantami są odpowiednio
$\|\partial_j\varphi\|_\infty |u|\chi_{B_R(0)}$
i $\|\varphi\|_\infty |v_j|\chi_{B_R(0)}$,
które należą do $L^1$ dzięki lokalnej całkowalności
(lub oszacowaniu Höldera na kuli).
Stąd $\partial_j u=v_j$ w sensie słabym i $u\in W^{1,p}$.

Odwrotnie, gdyby $u\in W^{1,p}$, to ograniczenie jej
słabych pochodnych do $\R^d\setminus\{0\}$ musiałoby
zgadzać się z klasycznymi pochodnymi $v_j$ prawie
wszędzie. Ponieważ singleton ma miarę zero, mielibyśmy
$v_j\in L^p$ dla wszystkich $j$, a więc również $|v|\in L^p$.
Obliczona całka radialna wymusza $(a+1)p<d$.
\end{przyklad}

\section{Zupełność i zamkniętość operatorów pochodnych}
\label{sec:pde9-complete}

\begin{twierdzenie}[Zupełność przestrzeni Sobolewa]
\label{thm:pde9-banach}
Dla każdego niepustego otwartego $\Omega\subset\R^d$, $k\in\N_0$ i
$1\le p\le\infty$ przestrzeń $W^{k,p}(\Omega)$ jest przestrzenią
Banacha. W szczególności $H^k(\Omega)$ jest przestrzenią Hilberta
z iloczynem skalarnym
\[
 \langle u,v\rangle_{H^k}=\sum_{|\alpha|\le k}
 \int_\Omega\partial^\alpha u(x)\,\overline{\partial^\alpha v(x)}
 \,\mathrm d\lambda_d(x).
\]
W przypadku zespolonym przyjmujemy liniowość względem
pierwszego argumentu; dla funkcji rzeczywistych sprzężenie
można pominąć.
\end{twierdzenie}
\begin{proof}
Niech $(u_n)$ będzie ciągiem Cauchy'ego w $W^{k,p}(\Omega)$.
Dla każdego $|\alpha|\le k$ zachodzi
$\|\partial^\alpha(u_n-u_m)\|_p\le\|u_n-u_m\|_{W^{k,p}}$,
także dla $p=\infty$. Zatem $(\partial^\alpha u_n)$
jest ciągiem Cauchy'ego w przestrzeni Banacha $L^p(\Omega)$.
Istnieje $f_\alpha\in L^p(\Omega)$ takie, że
$\partial^\alpha u_n\to f_\alpha$ w $L^p$. Oznaczmy $u=f_0$.
Dla dowolnego $\varphi\in C_c^\infty(\Omega)$ mamy
\[
 \int_\Omega u_n\partial^\alpha\varphi\,\mathrm d\lambda_d
 =(-1)^{|\alpha|}\int_\Omega
 (\partial^\alpha u_n)\varphi\,\mathrm d\lambda_d.
\]
Niech $p'$ będzie wykładnikiem sprzężonym do $p$,
z konwencją $1/\infty=0$. Test oraz każda jego pochodna
są ograniczone i mają zwarty nośnik, więc należą do $L^{p'}$,
również gdy $p'=1$ lub $p'=\infty$. Z nierówności Höldera
otrzymujemy
\[
 \begin{aligned}
 \left|\int_\Omega(u_n-u)\partial^\alpha\varphi
 \,\mathrm d\lambda_d\right|
 &\le\|u_n-u\|_p\|\partial^\alpha\varphi\|_{p'}\to0,\\
 \left|\int_\Omega(\partial^\alpha u_n-f_\alpha)\varphi
 \,\mathrm d\lambda_d\right|
 &\le\|\partial^\alpha u_n-f_\alpha\|_p\|\varphi\|_{p'}\to0.
 \end{aligned}
\]
Przejście do granicy daje dokładnie tożsamość testową dla
$\partial^\alpha u=f_\alpha$ w $\mathcal D'(\Omega)$.
Ponieważ wszystkie $f_\alpha$ należą do $L^p$, mamy
$u\in W^{k,p}$. Dla $p<\infty$ zachodzi
\[
 \|u_n-u\|_{W^{k,p}}^p
 =\sum_{|\alpha|\le k}\|\partial^\alpha u_n-f_\alpha\|_p^p\to0,
\]
gdyż suma jest skończona. Dla $p=\infty$ maksimum skończenie
wielu norm dążących do zera także dąży do zera.
Dowodzi to zupełności.

Dla $p=2$ każda całka w podanym iloczynie skalarnym
jest bezwzględnie zbieżna z nierówności Cauchy'ego--Schwarza.
Liniowość pochodnych i własności iloczynu skalarnego $L^2$
dają liniowość w pierwszym argumencie, antyliniowość
w drugim i symetrię hermitowską. Ponadto
$\langle u,u\rangle_{H^k}=\|u\|_{W^{k,2}}^2$,
a równość z zerem daje $u=0$ w $L^2$, ponieważ suma zawiera
składnik $\alpha=0$. Zupełność wykazana powyżej
dowodzi, że jest to przestrzeń Hilberta.
\end{proof}

\begin{wniosek}[Zamkniętość pochodnej]
\label{cor:pde9-closed}
Dla ustalonego wielowskaźnika $\alpha$ określmy operator
$D_\alpha u=\partial^\alpha u$ w przestrzeni $L^p(\Omega)$
na jego naturalnej dziedzinie
\[
 \mathcal D(D_\alpha)=
 \{u\in L^p(\Omega):\partial^\alpha u
 \text{ jest reprezentowana przez funkcję z }L^p(\Omega)\}.
\]
Operator $D_\alpha$ jest zamknięty: jeżeli
$u_n\in\mathcal D(D_\alpha)$, $u_n\to u$ w $L^p$ oraz
$D_\alpha u_n\to g$ w $L^p$, to $u\in\mathcal D(D_\alpha)$
i $D_\alpha u=g$.
Nie zakładamy istnienia ani zbieżności pozostałych
pochodnych rzędu $|\alpha|$ w $L^p$.
\end{wniosek}
\begin{proof}
Dla każdego testu $\varphi$ mamy
\[
 \int_\Omega u_n\partial^\alpha\varphi\,\mathrm d\lambda_d
 =(-1)^{|\alpha|}\int_\Omega(D_\alpha u_n)\varphi\,\mathrm d\lambda_d.
\]
Oszacowania Höldera z dowodu twierdzenia
~\ref{thm:pde9-banach}, zastosowane do $u_n-u$ oraz
$D_\alpha u_n-g$, pozwalają przejść do granicy po obu stronach.
Otrzymujemy $\partial^\alpha u=g$ w sensie dystrybucji,
co daje dokładnie żądaną przynależność do dziedziny i wartość
operatora. W normowanej przestrzeni $L^p\times L^p$
domkniętość wykresu jest równoważna tej własności ciągowej.
\end{proof}

\begin{uwaga}[Dziedzina pojedynczej pochodnej a przestrzeń Sobolewa]
Zamkniętość pojedynczej pochodnej nie oznacza, że granica
należy do całej przestrzeni $W^{|\alpha|,p}$.
Na przykład na $\Omega=(0,1)^2$ funkcja
$u(x,y)=\chi_{(1/2,1)}(y)$ należy do każdego $L^p(\Omega)$
i ma pochodną $\partial_x u=0$, ponieważ całka z
$\partial_x\varphi(x,y)$ po $x\in(0,1)$ jest zerowa dla każdego
$y$. Jednak $\partial_y u$ działa przez
\[
 \langle\partial_y u,\varphi\rangle
 =\int_0^1\varphi(x,1/2)\,\mathrm d\lambda_1(x).
\]
Nie jest to dystrybucja regularna. Istotnie, wybieramy
$\eta\in C_c^\infty(0,1)$ o całce $1$ i
$\psi\in C_c^\infty(-1,1)$ z $\psi(0)=1$.
Testy $\varphi_\varepsilon(x,y)=\eta(x)\psi((y-1/2)/\varepsilon)$
dla $0<\varepsilon<1/4$ mają jednostajnie ograniczone moduły,
wspólny zwarty nośnik i nośniki zawarte w pasach,
których miary dążą do zera. Działanie $\partial_y u$
na tych testach jest stale równe $1$, natomiast działanie
każdej funkcji lokalnie całkowalnej dążyłoby do zera
z absolutnej ciągłości całki Lebesgue'a.
Stąd $u\notin W^{1,p}(\Omega)$.
Dla $p<\infty$ można też aproksymować $u$ w $L^p$
gładkimi funkcjami zależnymi tylko od $y$ i mającymi
$\partial_x u_n=0$: wystarczy zastąpić skok gładkim
przejściem z $0$ do $1$ na przedziale
$(1/2-1/n,1/2+1/n)$, dla $n\ge3$.
Błąd w $L^p$ jest nie większy niż $(2/n)^{1/p}$.
To pokazuje, dlaczego we wniosku używamy naturalnej dziedziny,
a nie ograniczenia operatora do całej przestrzeni $W^{1,p}$.

Sama norma Sobolewa nie przypisuje klasie funkcji
kanonicznych wartości punktowych. Warunki istnienia
ciągłych reprezentantów wynikają z twierdzeń o zanurzeniach
Sobolewa, omawianych w następnym rozdziale.
\end{uwaga}


\section{Wymiar jeden: absolutna ciągłość}
\label{sec:pde9-oneD}

\begin{twierdzenie}[Charakterystyka $W^{1,p}(a,b)$]
\label{thm:pde9-ac}
Niech $-\infty<a<b<\infty$ i $1\le p\le\infty$.
Klasa $u\in L^p(a,b)$ należy do $W^{1,p}(a,b)$ wtedy i tylko
wtedy, gdy ma reprezentanta absolutnie ciągłego
$\widetilde u:[a,b]\to\R$, którego pochodna klasyczna
istnieje prawie wszędzie i należy do $L^p(a,b)$.
Reprezentant ten jest jedyny wśród funkcji ciągłych na $[a,b]$.
Jego pochodna klasyczna jest słabą pochodną $u'$ prawie wszędzie
i zachodzi
\[
 \widetilde u(y)-\widetilde u(x)
 =\int_x^y u'(t)\,\mathrm d\lambda_1(t)
 \qquad(a\le x\le y\le b).
\]
\end{twierdzenie}
\begin{proof}
Załóżmy $u\in W^{1,p}(a,b)$ i niech $g=u'$ będzie jego
słabą pochodną. Ponieważ przedział ma skończoną miarę $L=b-a$,
mamy $\|g\|_1\le L^{1-1/p}\|g\|_p$ dla $p<\infty$
oraz $\|g\|_1\le L\|g\|_\infty$ dla $p=\infty$.
Zdefiniujmy na $[a,b]$
\[
 U(x)=\int_a^x g(t)\,\mathrm d\lambda_1(t).
\]
Z podstawowego twierdzenia rachunku całkowego dla całki
Lebesgue'a wynika, że $U$ jest absolutnie ciągła i
$U'=g$ prawie wszędzie. Dla przypomnienia, absolutna
ciągłość wynika z oszacowania sumy przyrostów $U$
na parami rozłącznych przedziałach przez całkę z $|g|$
na ich sumie oraz z absolutnej ciągłości tej całki.
Dla testu $\varphi\in C_c^\infty(a,b)$ całkowanie przez
części dla funkcji absolutnie ciągłych daje
\[
 \int_a^b U(t)\varphi'(t)\,\mathrm d\lambda_1(t)
 =-\int_a^b g(t)\varphi(t)\,\mathrm d\lambda_1(t).
\]
Można zastosować je do iloczynu $U\varphi$, który jest
absolutnie ciągły, ma pochodną $g\varphi+U\varphi'$
prawie wszędzie i znika w obu końcach.
Stąd dystrybucyjna pochodna $u-U$ jest zerowa.
Twierdzenie~\ref{thm:pde7-constant} zastosowane na spójnym
przedziale $(a,b)$ daje $T_{u-U}=T_C$ dla pewnej stałej $C$.
Ponieważ $U$ jest ciągła na zwartym przedziale,
jest ograniczona, więc należy do $L^p(a,b)$.
Jednoznaczność reprezentanta dystrybucji regularnej
daje $u-U=C$ prawie wszędzie; dla funkcji rzeczywistej $C$
jest rzeczywista. Wobec tego $\widetilde u=U+C$ jest
żądanym reprezentantem i ma pochodną $g$.
Dla $a\le x\le y\le b$ addytywność całki daje podany wzór.

Odwrotnie, jeśli $u$ ma reprezentanta absolutnie ciągłego
$\widetilde u$ z pochodną $g\in L^p$, to iloczyn
$\widetilde u\varphi$ jest absolutnie ciągły, ma pochodną
$g\varphi+\widetilde u\varphi'$ prawie wszędzie
i znika w $a,b$. Całka z tej pochodnej jest zatem równa zeru.
Ponieważ $u=\widetilde u$ prawie wszędzie,
otrzymujemy $\int_a^b u\varphi'\,\mathrm d\lambda_1
=-\int_a^b g\varphi\,\mathrm d\lambda_1$.
To jest definicja słabej pochodnej z $L^p$.
Jeśli dwa ciągłe reprezentanty tej samej klasy różniłyby się
w dowolnym punkcie $[a,b]$, ich różnica byłaby niezerowa
na względnie otwartym otoczeniu tego punktu, którego przecięcie
z $(a,b)$ ma dodatnią miarę. Przeczy to równości prawie wszędzie.
\end{proof}

\begin{wniosek}[Oszacowanie punktowe i ślady w jednym wymiarze]
\label{cor:pde9-trace-oneD}
Dla $I=(a,b)$, $L=b-a$ i $u\in W^{1,p}(I)$,
$1\le p<\infty$, jedyny ciągły reprezentant spełnia
\[
 \|\widetilde u\|_{C([a,b])}\le
 L^{-1/p}\|u\|_{L^p(I)}+L^{1-1/p}\|u'\|_{L^p(I)}.
\]
Dla $p=\infty$ mamy
$\|\widetilde u\|_{C([a,b])}=\|u\|_\infty
\le\|u\|_{W^{1,\infty}}$.
W szczególności odwzorowania $u\mapsto\widetilde u(a)$
i $u\mapsto\widetilde u(b)$ są ciągłymi funkcjonałami
liniowymi na $W^{1,p}(I)$.
\end{wniosek}
\begin{proof}
Dla ciągłego reprezentanta supremum na $[a,b]$ jest równy
istotnemu supremum na $(a,b)$: dowolna wartość
ściśle mniejsza od maksimum modułu jest przekroczona
na zbiorze dodatniej miary dzięki ciągłości, także gdy
maksimum jest osiągnięte na końcu przedziału.
Dla $p<\infty$ wybierzmy $x_0\in[a,b]$,
w którym $|\widetilde u|$ osiąga minimum.
Mamy $L|\widetilde u(x_0)|^p\le\|u\|_p^p$.
Z twierdzenia~\ref{thm:pde9-ac}, dla dowolnego $x\in[a,b]$,
\[
 |\widetilde u(x)-\widetilde u(x_0)|
 \le\int_{\min(x,x_0)}^{\max(x,x_0)}|u'(t)|
 \,\mathrm d\lambda_1(t)
 \le L^{1-1/p}\|u'\|_p.
\]
Dodając oszacowanie minimum i biorąc supremum,
otrzymujemy tezę. Dla $p=\infty$ korzystamy z równości
supremum i istotnego supremum oraz definicji normy.
Jednoznaczność reprezentanta zapewnia liniowość
odwzorowań ewaluacji: reprezentant kombinacji liniowej
jest kombinacją liniową reprezentantów.
Dla $p<\infty$ moduł każdej ewaluacji jest ograniczony
przez $(L^{-1/p}+L^{1-1/p})\|u\|_{W^{1,p}}$,
a dla $p=\infty$ przez $\|u\|_{W^{1,\infty}}$.
\end{proof}

\begin{wniosek}[Moduł ciągłości i przypadek Lipschitza]
Dla $u\in W^{1,p}(a,b)$ i $1<p<\infty$ zachodzi
\[
 |\widetilde u(y)-\widetilde u(x)|
 \le\|u'\|_p|y-x|^{1-1/p}\qquad(x,y\in[a,b]).
\]
Dla $p=\infty$ reprezentant jest lipschitzowski ze stałą
co najwyżej $\|u'\|_\infty$. Odwrotnie, każda funkcja
lipschitzowska na $[a,b]$ wyznacza element $W^{1,\infty}(a,b)$.
Dla $p=1$ otrzymujemy absolutną ciągłość, ale nie wynika
stąd żaden ustalony dodatni wykładnik Höldera.
\end{wniosek}
\begin{proof}
Dla $x=y$ nierówność jest natychmiastowa; zamiana
$x,y$ nie zmienia modułu różnicy. Dla $x<y$
stosujemy do wzoru z twierdzenia
~\ref{thm:pde9-ac} nierówność Höldera na $(x,y)$:
$\int_x^y|u'|\,\mathrm d\lambda_1
\le |y-x|^{1-1/p}\|u'\|_{L^p(x,y)}$.
Dla $p=\infty$ całka jest ograniczona przez
$|y-x|\|u'\|_\infty$. Funkcja lipschitzowska jest absolutnie
ciągła, gdyż suma modułów jej przyrostów jest ograniczona
przez stałą Lipschitza razy sumę długości przedziałów.
Jej pochodna istnieje zatem prawie wszędzie
i jest ograniczona przez tę stałą z oszacowania ilorazów
różnicowych; sama funkcja jest ograniczona na $[a,b]$.
Charakterystyka z twierdzenia daje przynależność
do $W^{1,\infty}$. W szczególności najlepsza stała
Lipschitza reprezentanta jest równa $\|u'\|_\infty$.

Dla $p=1$ wzór całkowy daje kontrolę przez
$\int_x^y|u'|\,\mathrm d\lambda_1$, która dąży do zera
jednostajnie względem przedziałów o długości dążącej do zera.
Na przykład na $(0,1)$ dla dowolnego $\gamma>0$ wybierzmy
$0<\theta<\min\{1,\gamma\}$. Funkcja $u(x)=x^\theta$
jest absolutnie ciągła, gdyż jest całką z
$\theta t^{\theta-1}\in L^1(0,1)$ od $0$ do $x$.
Należy do $W^{1,1}(0,1)$, lecz
$|u(x)-u(0)|/x^\gamma=x^{\theta-\gamma}\to+\infty$
przy $x\downarrow0$. Nie ma więc wspólnego dodatniego
wykładnika Höldera dla całej przestrzeni $W^{1,1}$.
\end{proof}

\begin{przyklad}[Funkcja Cantora a pochodna słaba]
Klasyczna funkcja Cantora $F$ na $[0,1]$ jest ciągła,
niemalejąca, spełnia $F(0)=0$, $F(1)=1$ i jest stała na każdym
przedziale usuniętym w konstrukcji zbioru Cantora.
Zbiór Cantora ma miarę zero, ponieważ po $n$ krokach
jest zawarty w $2^n$ przedziałach o długości $3^{-n}$,
o łącznej długości $(2/3)^n$.
Wobec tego klasyczna pochodna $F'$ istnieje i jest równa zeru
poza zbiorem miary zero. Gdyby $F$ należała do $W^{1,1}(0,1)$,
twierdzenie~\ref{thm:pde9-ac} i jednoznaczność ciągłego
reprezentanta oznaczałyby, że jest absolutnie ciągła
i jej słaba pochodna jest zerowa prawie wszędzie.
Wzór całkowy dawałby wówczas $F(1)-F(0)=0$,
wbrew $F(1)-F(0)=1$.
Zatem pochodna istniejąca prawie wszędzie nie musi być
pochodną słabą. Ponieważ na $(0,1)$ zachodzi
$W^{1,p}\subset W^{1,1}$ dla wszystkich $1\le p\le\infty$
przez nierówność Höldera, $F$ nie należy do żadnej z tych
przestrzeni Sobolewa.
\end{przyklad}

\section{Sploty i aproksymacja wewnętrzna}
\label{sec:pde9-mollify}

Ustalmy nieujemną funkcję $\rho\in C_c^\infty(B_1(0))$
o całce $1$ i połóżmy
$\rho_\varepsilon(x)=\varepsilon^{-d}\rho(x/\varepsilon)$,
jak w sekcji~\ref{sec:pde7-mollifiers}.
Niech $\Omega_\varepsilon=\{x\in\Omega:
\operatorname{dist}(x,\R^d\setminus\Omega)>\varepsilon\}$,
przy czym odległość od zbioru pustego jest nieskończona.
Dla $u\in L^1_{\mathrm{loc}}(\Omega)$ i $x\in\Omega_\varepsilon$
definiujemy
\[
 u_\varepsilon(x)=(\rho_\varepsilon*u)(x)
 =\int_\Omega \rho_\varepsilon(x-y)u(y)\,\mathrm d\lambda_d(y).
\]
Całka ma sens, ponieważ jądro jest ograniczone
i jako funkcja zmiennej $y$ ma zwarty nośnik w $\Omega$.
Nie wymagamy przedłużenia $u$ przez zero na całej przestrzeni.

\begin{twierdzenie}[Lokalne wygładzanie Sobolewa]
\label{thm:pde9-mollify}
Niech $1\le p<\infty$, $k\in\N_0$ oraz
$u\in W^{k,p}_{\mathrm{loc}}(\Omega)$.
Funkcja $u_\varepsilon$ jest gładka na $\Omega_\varepsilon$
i dla każdego $|\alpha|\le k$ zachodzi
\[
 \partial^\alpha u_\varepsilon
 =\rho_\varepsilon*(\partial^\alpha u)
 \quad\text{na }\Omega_\varepsilon.
\]
Dla każdego zwartego $K\subset\Omega$ mamy dla
dostatecznie małych $\varepsilon>0$ zbieżność
\[
 \sum_{|\alpha|\le k}
 \|\partial^\alpha u_\varepsilon-\partial^\alpha u\|_{L^p(K)}
 \longrightarrow0\qquad(\varepsilon\downarrow0).
\]
W szczególności $u_\varepsilon\to u$ w $W^{k,p}(U)$
dla każdego otwartego $U\Subset\Omega$.
\end{twierdzenie}
\begin{proof}
Z sekcji~\ref{sec:pde9-def} mamy
$\partial^\alpha u=v_\alpha$ w sensie dystrybucji,
gdzie $v_\alpha\in L^p_{\mathrm{loc}}(\Omega)$,
włącznie z $v_0=u$. Każda z tych funkcji jest lokalnie
całkowalna. Twierdzenie~\ref{thm:pde7-mollify}
zastosowane do dystrybucji regularnej $u$ daje
gładkość i tożsamość pochodnych na $\Omega_\varepsilon$.
Obie strony są w tym miejscu gładkimi funkcjami.
Jeśli $K\ne\varnothing$ jest zwarty w $\Omega$,
to dla pewnego $r>0$ zbiór
$K+\overline{B_{2r}(0)}$ jest zwarty w $\Omega$.
W szczególności $K\subset\Omega_\varepsilon$ dla
$0<\varepsilon<r$ (dla $K=\varnothing$ teza jest pusta).
Z twierdzenia~\ref{thm:pde7-approx-id}, zastosowanego
osobno do każdej funkcji $v_\alpha$, wynika
\[
 \|\partial^\alpha u_\varepsilon-v_\alpha\|_{L^p(K)}
 =\|\rho_\varepsilon*v_\alpha-v_\alpha\|_{L^p(K)}
 \longrightarrow0.
\]
Ponieważ wielowskaźników jest skończenie wiele,
ich suma dąży do zera. Dla $K=\overline U$ ograniczamy
każdą normę do $U$ i stosujemy równoważność norm
z sekcji~\ref{sec:pde9-def}.
Nie definiujemy tutaj przestrzeni Sobolewa na dowolnym
zbiorze zwartym: zapis dotyczy ograniczeń pochodnych
z otoczenia $K$ do tego zbioru.
\end{proof}

\begin{lemat}[Odcięcie i wygładzanie o zwartym nośniku]
\label{lem:pde9-compact-product}
Niech $u\in W^{k,p}_{\mathrm{loc}}(\Omega)$,
$1\le p\le\infty$, oraz $\psi\in C_c^\infty(\Omega)$.
Przedłużenie $w$ funkcji $\psi u$ przez zero poza $\Omega$
należy do $W^{k,p}(\R^d)$ i dla $|\alpha|\le k$
jego słaba pochodna jest przedłużeniem przez zero funkcji
\[
 \partial^\alpha(\psi u)=
 \sum_{\beta\le\alpha}\binom{\alpha}{\beta}
 (\partial^\beta\psi)(\partial^{\alpha-\beta}u).
\]
Jeżeli ponadto $p<\infty$, to
$\rho_\varepsilon*w\to w$ w $W^{k,p}(\R^d)$.
Dla dowolnego otwartego $V\Subset\Omega$
zawierającego $\supp\psi$ i dostatecznie małego
$\varepsilon>0$ mamy $\rho_\varepsilon*w\in C_c^\infty(V)$.
\end{lemat}
\begin{proof}
Dla $\varphi\in C_c^\infty(\R^d)$ funkcja
$\psi\varphi|_\Omega$ jest testem na $\Omega$.
W pierwszym rzędzie, dla $k\ge1$, otrzymujemy
\[
 \begin{aligned}
 \langle\partial_j w,\varphi\rangle
 &=-\int_\Omega u\psi\,\partial_j\varphi\,\mathrm d\lambda_d\\
 &=-\int_\Omega u\,\partial_j(\psi\varphi)\,\mathrm d\lambda_d
   +\int_\Omega u(\partial_j\psi)\varphi\,\mathrm d\lambda_d\\
 &=\int_\Omega(\psi\partial_j u+u\partial_j\psi)
 \varphi\,\mathrm d\lambda_d.
 \end{aligned}
\]
Wszystkie funkcje po prawej stronie mają nośniki
zawarte w $\supp\psi$. Ta tożsamość dotyczy również
gładkich czynników $\partial^\beta\psi$ oraz słabych
pochodnych $u$ rzędu mniejszego niż $k$.
Iterując ją i korzystając ze wzoru Pascala dla
współczynników dwumianowych, otrzymujemy podany
wzór wielowskaźnikowy. W każdym składniku gładki
czynnik jest ograniczony i ma zwarty nośnik w $\Omega$,
na którego otoczeniu drugi czynnik należy do $L^p$.
Zatem wszystkie przedłużone pochodne należą do
$L^p(\R^d)$, także dla $p=\infty$.
Dla $k=0$ wystarcza ten sam argument dla $w$.
Nie korzystamy z twierdzenia o przedłużaniu przez zero
dowolnego elementu przestrzeni $W^{k,p}_0$.

Niech $A=\supp\psi$. Jeśli $A=\varnothing$, to $w=0$.
W przeciwnym przypadku $A$ jest zwarty, a nośniki
$w$ i wszystkich wskazanych pochodnych są zawarte w $A$
w sensie zanikania prawie wszędzie poza $A$.
Z definicji splotu dla $0<\varepsilon\le1$
wynika, że $\rho_\varepsilon*w$ znika poza
$A+\overline{B_1(0)}$. To samo dotyczy jego pochodnych,
gdyż funkcja jest gładka. Stosujemy twierdzenie
~\ref{thm:pde9-mollify} na całej $\R^d$ ze wspólnym
zbiorem zwartym $A+\overline{B_1(0)}$.
Dla $p<\infty$ zbieżność lokalna norm staje się
globalną, ponieważ różnice pochodnych znikają poza
tym zbiorem. Wreszcie
$A+\overline{B_\varepsilon(0)}\subset V$ dla
$\varepsilon<\operatorname{dist}(A,\R^d\setminus V)$,
co daje wymagany zwarty nośnik w $V$.
\end{proof}

\begin{uwaga}[Dlaczego nie $p=\infty$?]
Konkluzja normowa na ogół zawodzi dla $p=\infty$,
nawet dla $k=1$. Na $I=(-1,1)$ funkcja $u(x)=|x|$
ma słabą pochodną $\operatorname{sgn}x$.
Dla dowolnego $0<r<1$ i każdej funkcji
$v\in C^1((-r,r))$ zachodzi
\[
 \|v'-\operatorname{sgn}\|_{L^\infty(-r,r)}\ge1.
\]
Jeżeli norma po lewej stronie jest skończona i równa $M$,
to na obu półprzedziałach nierówności
$|v'(x)+1|\le M$ dla $x<0$ i $|v'(x)-1|\le M$ dla $x>0$
zachodzą w każdym punkcie dzięki ciągłości i temu,
że każdy niepusty przedział otwarty ma dodatnią miarę.
Przechodząc do zera z obu stron, dostajemy
$2\le|v'(0)+1|+|v'(0)-1|\le2M$.
W szczególności dotyczy to każdego wygładzenia $u_\varepsilon$,
jeśli jest określone na $(-r,r)$.
Jednak dla $u\in W^{k,\infty}_{\mathrm{loc}}(\Omega)$
i $1\le q<\infty$ mamy
$u\in W^{k,q}_{\mathrm{loc}}(\Omega)$: na każdym
$U\Subset\Omega$ stosujemy oszacowanie
$\|v_\alpha\|_{L^q(U)}\le
\lambda_d(U)^{1/q}\|v_\alpha\|_{L^\infty(U)}$.
Twierdzenie~\ref{thm:pde9-mollify} daje więc lokalną
zbieżność w każdej takiej skończonej normie.
\end{uwaga}

\begin{center}
\begin{tikzpicture}[scale=.92,>=Stealth]
 \draw[thick,rounded corners=12pt,fill=blue!5,draw=blue!55]
  (-4.1,-1.55) -- (3.9,-1.55) -- (4.2,1.55) -- (-3.7,1.65) -- cycle;
 \draw[thick,rounded corners=10pt,fill=green!9,draw=green!48!black]
  (-2.9,-1.15) rectangle (2.9,1.15);
 \draw[thick,rounded corners=9pt,fill=orange!11,draw=orange!75!black]
  (-1.55,-.65) rectangle (1.55,.65);
 \node at (-3.45,1.18) {$\Omega$};
 \node at (2.48,.86) {$U$};
 \node at (-1.05,.30) {$K$};
 \coordinate (x) at (.95,-.05);
 \draw[blue!65!black,thick,dashed] (x) circle (.8);
 \fill (x) circle (1.5pt) node[below left] {$x$};
 \draw[->,blue!65!black,thick] (x)--(1.75,-.05)
  node[midway,above] {$\varepsilon$};
 \node[blue!65!black,anchor=west] at (1.82,-.38) {$B_\varepsilon(x)$};
\end{tikzpicture}

{\small\emph{Lokalność wygładzania.} Dla $K\Subset U\Subset\Omega$
wybieramy $\varepsilon<\operatorname{dist}(K,\R^d\setminus U)$.
Wówczas dla każdego $x\in K$ nośnik jądra
$y\mapsto\rho_\varepsilon(x-y)$ jest zawarty w $U$.
Splot korzysta wyłącznie z wartości funkcji w tym otoczeniu.}
\end{center}


\section{Twierdzenie Meyersa--Serrina i gęstość funkcji gładkich}
\label{sec:pde9-density}

\begin{twierdzenie}[Meyers--Serrin]
\label{thm:pde9-MeyersSerrin}
Niech $\Omega\subset\R^d$ będzie niepustym zbiorem otwartym,
$k\in\N_0$ i $1\le p<\infty$. Wówczas
$C^\infty(\Omega)\cap W^{k,p}(\Omega)$
jest gęsta w $W^{k,p}(\Omega)$.
Nie zakłada się regularności brzegu $\partial\Omega$.
\end{twierdzenie}
\begin{proof}
Najpierw skonstruujemy potrzebny rozkład jedynki
z kontrolą nośników. Dla $j\ge1$ połóżmy
\[
 K_j=\{x\in\Omega:|x|\le j,\
 \operatorname{dist}(x,\R^d\setminus\Omega)\ge1/j\},
\]
a dla $j\le0$ przyjmijmy $K_j=\varnothing$.
Dla $\Omega=\R^d$ drugi warunek pomijamy.
Zbiory $K_j$ są zwarte, gdyż są ograniczone
i domknięte w $\R^d$ (dla niepustego dopełnienia
odległość jest funkcją ciągłą, a dodatnia odległość
wyklucza punkty z dopełnienia).
Ścisłe nierówności $j<j+1$ oraz $1/j>1/(j+1)$
dają $K_j\subset\operatorname{int}K_{j+1}$,
a każdy punkt $\Omega$ należy do wnętrza pewnego $K_j$.
W szczególności każdy zwarty podzbiór $\Omega$
zawiera się we wnętrzu pewnego $K_j$, przez skończone
podpokrycie i monotoniczność zbiorów.

Dla $j\ge1$ wybierzmy $\eta_j\in C_c^\infty(\operatorname{int}K_{j+1})$
z $0\le\eta_j\le1$ i $\eta_j=1$ w otoczeniu $K_j$.
Jeśli $K_j=\varnothing$, przyjmujemy $\eta_j=0$.
W pozostałych przypadkach konstrukcja jest następująca:
wybieramy $r>0$, dla którego
$K_j+\overline{B_{3r}(0)}\subset\operatorname{int}K_{j+1}$.
Splot
$\eta_j=\rho_r*\chi_{K_j+\overline{B_{2r}(0)}}$
jest gładki dzięki twierdzeniu~\ref{thm:pde7-mollify}
na $\R^d$, przyjmuje wartości w $[0,1]$,
jest równy $1$ na $K_j+B_r(0)$ i ma nośnik
w $K_j+\overline{B_{3r}(0)}$.
Wynika to z nieujemności, normalizacji i nośnika jądra.
W szczególności $\eta_{j+1}=1$ na $\supp\eta_j$,
więc $\eta_j\le\eta_{j+1}$.
Połóżmy $\eta_0=0$ i $\psi_j=\eta_j-\eta_{j-1}$.
Funkcje $\psi_j$ są nieujemne, gładkie i mają zwarty nośnik
w $\Omega$, a przez teleskopowanie
$\sum_{j=1}^N\psi_j=\eta_N$.
Dla każdego punktu $x\in\Omega$ mamy $\eta_N(x)=1$
od pewnego $N$, stąd $\sum_{j\ge1}\psi_j=1$.

Ustalmy otwarte zbiory
$V_j=\operatorname{int}K_{j+2}\setminus K_{j-2}$,
$ j\ge1$. Mamy $V_j\Subset\Omega$, a
$\supp\psi_j\subset V_j$. Istotnie, oba składniki
$\psi_j$ mają nośniki w $\operatorname{int}K_{j+1}$
i dla $j\ge2$ oba są równe $1$ w otoczeniu $K_{j-1}$,
jeśli zbiór ten jest niepusty. Dla małych $j$
zbiory o niedodatnich indeksach są puste.
Rodzina $V_j$ jest lokalnie skończona:
jeśli zwarty $A\subset\Omega$ zawiera się
we wnętrzu $K_N$, to $A\cap V_j=\varnothing$
dla $j\ge N+2$. Zatem także rodzina nośników $\psi_j$
jest lokalnie skończona. Nie wymagaliśmy skończonej
liczby zbiorów ani jakiejkolwiek regularności brzegu.

Niech teraz $u\in W^{k,p}(\Omega)$ i $\delta>0$.
Dla każdego $j$ oznaczmy przez $w_j$ przedłużenie
$\psi_j u$ przez zero na $\R^d$.
Lemat~\ref{lem:pde9-compact-product} daje
$w_j\in W^{k,p}(\R^d)$ oraz pozwala wybrać
$\varepsilon_j>0$ tak, aby $g_j=\rho_{\varepsilon_j}*w_j$
należało do $C_c^\infty(V_j)$ i
\[
 \|g_j-w_j\|_{W^{k,p}(\R^d)}<\delta2^{-j}.
\]
Dla $\psi_j=0$ bierzemy $g_j=0$.
Ograniczenie normy do $\Omega$ nie zwiększa błędu.
Funkcja $g=\sum_{j\ge1}g_j|_\Omega$ jest gładka,
ponieważ lokalnie suma ma tylko skończenie wiele
niezerowych składników. Nie zakładamy jeszcze,
że $g\in W^{k,p}(\Omega)$.

Niech $e_j=g_j|_\Omega-w_j|_\Omega$.
Szereg $\sum_{j\ge1}e_j$ zbiega w przestrzeni Banacha
$W^{k,p}(\Omega)$ z twierdzenia~\ref{thm:pde9-banach}:
dla $M>N$ norma sumy od $N+1$ do $M$ nie przekracza
$\delta\sum_{j>N}2^{-j}$, co dąży do zera.
Oznaczmy jego granicę przez $e$; ciągłość normy daje
$\|e\|_{W^{k,p}}\le\delta$.
Wykażemy, że $e=g-u$ prawie wszędzie.
Dla dowolnego $\varphi\in C_c^\infty(\Omega)$
zbieżność w $L^p$ oraz nierówność Höldera dają
\[
 \int_\Omega e\varphi\,\mathrm d\lambda_d
 =\lim_{N\to\infty}\sum_{j=1}^N
 \int_\Omega e_j\varphi\,\mathrm d\lambda_d.
\]
Tylko skończenie wiele składników może mieć
nośnik przecinający $\supp\varphi$.
Na otoczeniu tego nośnika mamy
$\sum_j w_j|_\Omega=u$ prawie wszędzie,
ponieważ $\sum_j\psi_j=1$.
Prawa strona jest zatem równa
$\int_\Omega(g-u)\varphi\,\mathrm d\lambda_d$.
Funkcja $g-u$ jest lokalnie całkowalna, więc
jednoznaczność dystrybucji regularnych daje
$e=g-u$ prawie wszędzie. W konsekwencji
$g=u+e\in W^{k,p}(\Omega)$ oraz
$\|g-u\|_{W^{k,p}}\le\delta$.
Dla każdego zadanego błędu $\epsilon>0$
wybieramy $\delta=\epsilon/2$, otrzymując tezę.
Nie musimy zakładać zbieżności szeregu $\sum_j w_j$
w globalnej normie Sobolewa; zbieżny w tej normie
jest szereg błędów aproksymacji.
\end{proof}

\begin{uwaga}[Gładkość nie oznacza zerowania na brzegu]
Twierdzenie~\ref{thm:pde9-MeyersSerrin} nie głosi,
że $C_c^\infty(\Omega)$ jest gęsta w $W^{k,p}(\Omega)$.
Funkcje o zwartym nośniku wewnątrz $\Omega$
znikają w pewnym otoczeniu brzegu, podczas gdy
$C^\infty(\Omega)$ nie narzuca takiego warunku.
Gdy $\Omega=\R^d$ i $p<\infty$, gęstość
$C_c^\infty(\R^d)$ rzeczywiście zachodzi.
Aby to sprawdzić, wybierzmy $\psi\in C_c^\infty(\R^d)$,
$0\le\psi\le1$, równą $1$ na $B_1(0)$ i zerową
poza $B_2(0)$, jak w dowodzie twierdzenia
~\ref{thm:pde8-density}, oraz $\psi_R(x)=\psi(x/R)$.
Dla $u\in W^{k,p}(\R^d)$, $R\ge1$, z lematu
~\ref{lem:pde9-compact-product} wynika
$\psi_Ru\in W^{k,p}(\R^d)$.
W rozwinięciu
\[
 \partial^\alpha((1-\psi_R)u)=
 (1-\psi_R)\partial^\alpha u+
 \sum_{0<\beta\le\alpha}\binom{\alpha}{\beta}
 \partial^\beta(1-\psi_R)\partial^{\alpha-\beta}u
\]
pierwszy składnik ma normę $L^p$ ograniczoną przez
$\|\partial^\alpha u\|_{L^p(\R^d\setminus B_R(0))}\to0$.
W pozostałych składnikach gładki czynnik ma normę
supremum równą
$R^{-|\beta|}\|\partial^\beta\psi\|_\infty$,
a drugi ma skończoną normę $L^p$.
Każdy taki iloczyn dąży więc do zera w $L^p$.
Wzór wynika z liniowości pochodnych i lematu
~\ref{lem:pde9-compact-product}.
Otrzymujemy $\psi_Ru\to u$ w $W^{k,p}(\R^d)$.
Dla ustalonego $R$ stosujemy ponownie ten lemat,
aby aproksymować $\psi_Ru$ przez gładkie funkcje
o zwartym nośniku. Wybierając najpierw $R$,
a potem promień wygładzania z błędami mniejszymi
niż połowa zadanego błędu, dostajemy żądaną gęstość.

Dla $p=\infty$ nawet gęstość
$C^\infty(\Omega)\cap W^{k,\infty}(\Omega)$
nie zachodzi ogólnie: dla $k=1$ i $\Omega=(-1,1)$
dowiedziony w sekcji~\ref{sec:pde9-mollify}
dolny próg $1$ dla błędu pochodnej funkcji $|x|$
wyklucza taką aproksymację.
\end{uwaga}

\begin{przyklad}[Stała na przedziale nie należy do $W^{1,p}_0$]
\label{ex:pde9-constant}
Niech $\Omega=(0,1)$, $1\le p<\infty$, oraz $u\equiv1$.
Ponieważ $u\in L^p(0,1)$ i jej słaba pochodna
jest zerowa, $u\in W^{1,p}(0,1)$, ale
$u\notin W^{1,p}_0(0,1)$. Gdyby istniał ciąg
$\varphi_n\in C_c^\infty(0,1)$ zbieżny do $u$
w $W^{1,p}$, to z ciągłości śladu
(wniosek~\ref{cor:pde9-trace-oneD}) wynikałoby
\[
 |\varphi_n(0)-1|\le
 \|\varphi_n-u\|_{L^p(0,1)}+
 \|\varphi_n'-u'\|_{L^p(0,1)}\longrightarrow0.
\]
Jest to niemożliwe, ponieważ ciągły reprezentant
każdej $\varphi_n$ jest równy zero w $0$.
\end{przyklad}

\section[Domknięcie funkcji testowych i przedłużenie przez zero]{Przestrzeń $W^{k,p}_0$ i przedłużenie przez zero}
\label{sec:pde9-zero}

\begin{twierdzenie}[Przedłużenie przez zero]
\label{thm:pde9-zero-extension}
Niech $\Omega\subset\R^d$ będzie niepustym zbiorem otwartym,
$k\in\N_0$, $1\le p<\infty$ i $u\in W^{k,p}_0(\Omega)$.
Przedłużenie przez zero
\[
 E_0u(x)=\begin{cases}u(x),&x\in\Omega,\\0,&x\notin\Omega\end{cases}
\]
należy do $W^{k,p}(\R^d)$ i spełnia
\[
 \partial^\alpha(E_0u)=E_0(\partial^\alpha u)
 \quad (|\alpha|\le k),\qquad
 \|E_0u\|_{W^{k,p}(\R^d)}=\|u\|_{W^{k,p}(\Omega)}.
\]
Operator $E_0$ jest liniowym izometrycznym zanurzeniem
o domkniętym obrazie.
\end{twierdzenie}
\begin{proof}
Przedłużenie przez zero klas funkcji z $L^p(\Omega)$
jest dobrze określoną operacją liniową i zachowuje normę,
gdyż $\int_{\R^d}|E_0f|^p\,\mathrm d\lambda_d
=\int_\Omega|f|^p\,\mathrm d\lambda_d$.
Nie wynika z tego jeszcze zachowanie słabych pochodnych.

Z definicji $W^{k,p}_0$ istnieje ciąg
$\varphi_n\in C_c^\infty(\Omega)$ taki, że
$\varphi_n\to u$ w $W^{k,p}(\Omega)$.
Dla każdej $\varphi_n$ jej nośnik jest zwartym podzbiorem
$\Omega$, a zatem w otoczeniu każdego punktu brzegu
funkcja i jej przedłużenie przez zero są zerowe.
Stąd $E_0\varphi_n\in C_c^\infty(\R^d)$ i
$\partial^\alpha E_0\varphi_n=E_0(\partial^\alpha\varphi_n)$.
W szczególności dla $n,m$ mamy
\[
 \|E_0\varphi_n-E_0\varphi_m\|_{W^{k,p}(\R^d)}
 =\|\varphi_n-\varphi_m\|_{W^{k,p}(\Omega)}.
\]
Ciąg po lewej stronie jest więc ciągiem Cauchy'ego.
Z twierdzenia~\ref{thm:pde9-banach} istnieje jego
granica $v\in W^{k,p}(\R^d)$.
Jednocześnie
\[
 \|E_0\varphi_n-E_0u\|_{L^p(\R^d)}
 =\|\varphi_n-u\|_{L^p(\Omega)}\to0.
\]
Z jednoznaczności granicy w $L^p$ dostajemy $v=E_0u$.
Dla dowolnego $|\alpha|\le k$ analogicznie
\[
 \|\partial^\alpha E_0\varphi_n-E_0(\partial^\alpha u)\|_{L^p(\R^d)}
 =\|\partial^\alpha\varphi_n-\partial^\alpha u\|_{L^p(\Omega)}
 \to0.
\]
Ponieważ zbieżność w normie Sobolewa daje zarazem
$\partial^\alpha E_0\varphi_n\to\partial^\alpha v$ w $L^p$,
otrzymujemy wzór na pochodną. Sumując $p$-te potęgi norm
po wszystkich wielowskaźnikach, dostajemy równość norm.

Przestrzeń $W^{k,p}_0(\Omega)$ jest domkniętą podprzestrzenią
przestrzeni Banacha $W^{k,p}(\Omega)$, więc jest zupełna.
Jeżeli $E_0u_n$ zbiega w $W^{k,p}(\R^d)$, to z izometrii
$(u_n)$ jest ciągiem Cauchy'ego w $W^{k,p}_0(\Omega)$.
Dla jego granicy $u$ izometria daje $E_0u_n\to E_0u$.
Obraz operatora jest zatem domknięty.
\end{proof}

\begin{uwaga}[Znaczenie warunku brzegowego]
Nie każdą funkcję z przestrzeni Sobolewa można tak przedłużyć.
Dla $u\equiv1$ na $(0,1)$ mamy $u\in W^{1,p}(0,1)$,
lecz dla $\varphi\in C_c^\infty(\R)$ zachodzi
\[
 \langle (E_0u)',\varphi\rangle
 =-\int_0^1\varphi'(x)\,\mathrm d\lambda_1(x)
 =\varphi(0)-\varphi(1).
\]
Pochodną dystrybucyjną jest więc $\delta_0-\delta_1$.
Nie może jej reprezentować funkcja lokalnie całkowalna:
ograniczając testy do otoczenia zera niezawierającego $1$,
otrzymalibyśmy regularną reprezentację $\delta_0$,
co wyklucza argument z sekcji~\ref{sec:pde7-def}.
W konsekwencji $E_0u\notin W^{1,p}(\R)$.

Na dowolnym zbiorze otwartym nie zachodzi również
implikacja odwrotna do twierdzenia~\ref{thm:pde9-zero-extension}.
Niech $\Omega=(-1,0)\cup(0,1)$ oraz
$\zeta\in C_c^\infty(-1,1)$, $\zeta(0)=1$.
Dla $u=\zeta|_\Omega$ i $1\le p<\infty$ mamy
$u\in W^{1,p}(\Omega)$, a $E_0u=\zeta$ prawie wszędzie
na $\R$, ponieważ różnica może występować tylko w zerze.
Zatem $E_0u\in W^{1,p}(\R)$.
Gdyby $u\in W^{1,p}_0(\Omega)$, istniałyby
$\varphi_n\in C_c^\infty(\Omega)$ zbieżne do $u$
w $W^{1,p}(\Omega)$. Po ograniczeniu do $(0,1)$
zbieżność zachodziłaby nadal w $W^{1,p}(0,1)$.
Każda $\varphi_n$ ma jednak zerową wartość brzegową
w $0$ od prawej strony, podczas gdy $u$ ma tam wartość $1$.
Przeczy to ciągłości śladu z wniosku
~\ref{cor:pde9-trace-oneD}.
W szczególności samo przedłużenie przez zero należące
do $W^{1,p}(\R^d)$ nie charakteryzuje $W^{1,p}_0$
bez dodatkowych założeń o zbiorze otwartym.
\end{uwaga}


\section{Nierówność Poincarégo na kostce}
\label{sec:pde9-poincare}

\begin{twierdzenie}[Poincaré--Wirtinger na kostce]
\label{thm:pde9-cube-Poincare}
Niech $Q=\prod_{j=1}^d(a_j,a_j+L)$, $L>0$,
$1\le p\le\infty$. Dla $u\in W^{1,p}(Q)$ połóżmy
\[
 u_Q=\frac1{\lambda_d(Q)}\int_Q u\,\mathrm d\lambda_d.
\]
Średnia jest dobrze określona, ponieważ $Q$ ma skończoną
miarę i $L^p(Q)\subset L^1(Q)$.
Wówczas
\[
 \|u-u_Q\|_{L^p(Q)}
 \le L\sum_{j=1}^d\|\partial_j u\|_{L^p(Q)}.
\]
Prawa strona jest równa zeru dokładnie wtedy,
gdy $u$ jest prawie wszędzie stała na $Q$.
\end{twierdzenie}
\begin{proof}
Najpierw ustalmy $1\le p<\infty$ i niech $u$
będzie gładka w otoczeniu $\overline Q$.
Dla $h\in C^1([a,a+L])$ podstawowe twierdzenie
rachunku całkowego daje
\[
 \begin{aligned}
 \left|h(x)-\frac1L\int_a^{a+L}h(t)\,\mathrm d\lambda_1(t)\right|
 &\le\frac1L\int_a^{a+L}|h(x)-h(t)|\,\mathrm d\lambda_1(t)\\
 &\le\int_a^{a+L}|h'(s)|\,\mathrm d\lambda_1(s).
 \end{aligned}
\]
Ostatnia nierówność wynika stąd, że dla każdego $t$
moduł różnicy jest nie większy od całki z $|h'|$
na całym przedziale. Biorąc normę $L^p$ względem $x$
i stosując nierówność Höldera, otrzymujemy
\[
 \left\|h-\frac1L\int_a^{a+L}h\,\mathrm d\lambda_1\right\|_p
 \le L^{1/p}\|h'\|_1
 \le L\|h'\|_p.
\]

Dla $j=1,\ldots,d$ niech $I_j=(a_j,a_j+L)$ i
niech $A_j$ oznacza uśrednienie względem pierwszych
$j$ współrzędnych:
\[
 (A_j u)(x)=\frac1{L^j}
 \int_{\prod_{i=1}^j I_i}u(t_1,\ldots,t_j,x_{j+1},\ldots,x_d)
 \,\mathrm d\lambda_j(t).
\]
Traktujemy $A_j u$ jako funkcję na całej kostce,
stałą w pierwszych $j$ zmiennych. Przyjmujemy $A_0u=u$;
wówczas $A_du=u_Q$. Dla $v\in L^p(Q)$
nierówność Jensena względem znormalizowanej miary
na $\prod_{i=1}^j I_i$ daje $|A_jv|^p\le A_j(|v|^p)$
prawie wszędzie. Całkując po $Q$ i stosując Fubiniego,
otrzymujemy
\[
 \|A_jv\|_p^p\le\int_Q A_j(|v|^p)\,\mathrm d\lambda_d
 =\int_Q|v|^p\,\mathrm d\lambda_d=\|v\|_p^p.
\]
Zatem $A_j$ jest kontrakcją.

Dla ustalonych pozostałych współrzędnych stosujemy
nierówność jednowymiarową do funkcji $A_{j-1}u$
zmiennej $x_j$, której średnia w tej zmiennej jest
$A_ju$. Po podniesieniu do potęgi $p$ i scałkowaniu
względem pozostałych współrzędnych (dla $d=1$
nie ma drugiego całkowania) dostajemy
\[
 \|A_{j-1}u-A_ju\|_p^p
 \le L^p\|\partial_j A_{j-1}u\|_p^p.
\]
Różniczkowanie w $x_j$ komutuje z $A_{j-1}$,
ponieważ całkujemy względem innych współrzędnych,
a odpowiednie pochodne $u$ są ciągłe i ograniczone
na $\overline Q$. Z kontrakcyjności otrzymujemy
\[
 \|\partial_j A_{j-1}u\|_p
 =\|A_{j-1}(\partial_j u)\|_p\le\|\partial_j u\|_p.
\]
Teleskopowanie $u-u_Q=\sum_{j=1}^d(A_{j-1}u-A_ju)$
oraz nierówność trójkąta dają tezę dla gładkich $u$.

Niech teraz $u\in W^{1,p}(Q)$, nadal dla $p<\infty$.
Dla $0<\delta<L/2$ definiujemy kostkę
$Q_\delta=\prod_j(a_j+\delta,a_j+L-\delta)$.
Z twierdzenia~\ref{thm:pde9-mollify} wynika, że
dla ustalonego $\delta$ funkcje $u_\varepsilon$ są
gładkie w otoczeniu $\overline{Q_\delta}$ dla
dostatecznie małych $\varepsilon$ i zbiegają do $u$
w $W^{1,p}(Q_\delta)$. Ponadto
\[
 |(u_\varepsilon)_{Q_\delta}-u_{Q_\delta}|
 \le\lambda_d(Q_\delta)^{-1/p}
      \|u_\varepsilon-u\|_{L^p(Q_\delta)}\to0
\]
z nierówności Höldera. W szczególności
$u_\varepsilon-(u_\varepsilon)_{Q_\delta}\to
u-u_{Q_\delta}$ w $L^p(Q_\delta)$.
Stosujemy już udowodnione oszacowanie na $Q_\delta$
i przechodzimy z $\varepsilon$ do zera:
\[
 \|u-u_{Q_\delta}\|_{L^p(Q_\delta)}
 \le (L-2\delta)\sum_{j=1}^d\|\partial_j u\|_{L^p(Q_\delta)}.
\]
Teraz $u_{Q_\delta}\to u_Q$ przy $\delta\downarrow0$,
ponieważ $\lambda_d(Q_\delta)\to\lambda_d(Q)>0$
i całki z $u\in L^1(Q)$ zbiegają przez twierdzenie
o zbieżności zmajoryzowanej. Mamy też
\[
 \left|\|u-u_{Q_\delta}\|_{L^p(Q_\delta)}
       -\|u-u_Q\|_{L^p(Q_\delta)}\right|
 \le\lambda_d(Q_\delta)^{1/p}|u_{Q_\delta}-u_Q|\to0.
\]
Normy funkcji $u-u_Q$ oraz $\partial_j u$ na
$Q_\delta$ zbiegają do norm na $Q$ dzięki całkowalności
ich $p$-tych potęg i zbieżności zmajoryzowanej.
Przejście do granicy daje żądaną nierówność.

W przypadku $p=\infty$ każda funkcja
$u\in W^{1,\infty}(Q)$ należy również do $W^{1,q}(Q)$
dla każdego $1\le q<\infty$, z tymi samymi słabymi pochodnymi.
Dla dowolnego $f\in L^\infty(Q)$ mamy
$\|f\|_q\to\|f\|_\infty$ przy $q\to\infty$.
Istotnie, jeśli $M=\|f\|_\infty$, to
$\|f\|_q\le M\lambda_d(Q)^{1/q}$.
Dla $M>0$ i $0<\gamma<M$ zbiór
$A=\{x\in Q:|f(x)|>M-\gamma\}$ ma dodatnią miarę,
więc $\|f\|_q\ge(M-\gamma)\lambda_d(A)^{1/q}$.
Przejście z $q$ do nieskończoności, a następnie
z $\gamma$ do zera, daje równość granicy; dla $M=0$
jest ona oczywista. Stosujemy zatem oszacowanie dla
skończonego $q$ do $u$, a powyższą granicę do
$u-u_Q$ i każdej $\partial_j u$.

Jeśli wszystkie pierwsze pochodne są zerowe,
nierówność daje $\|u-u_Q\|_p=0$, czyli $u=u_Q$
prawie wszędzie. Odwrotnie, funkcja prawie wszędzie
stała ma zerowe pochodne dystrybucyjne, ponieważ
całka z pochodnej testu o zwartym nośniku jest zerowa.
\end{proof}

\begin{wniosek}[Wersja z zerowym warunkiem brzegowym]
\label{cor:pde9-Poincare-zero}
Dla $1\le p<\infty$ i $u\in W^{1,p}_0(Q)$
na kostce z poprzedniego twierdzenia zachodzi
\[
 \|u\|_{L^p(Q)}\le L\|\partial_j u\|_{L^p(Q)}
 \qquad (j=1,\ldots,d).
\]
W szczególności
\[
 |u|_{1,p}=\left(\sum_{j=1}^d\|\partial_j u\|_p^p\right)^{1/p}
\]
jest normą równoważną normie $W^{1,p}$ na $W^{1,p}_0(Q)$,
z oszacowaniami
$|u|_{1,p}\le\|u\|_{W^{1,p}}\le(1+L^p)^{1/p}|u|_{1,p}$.
\end{wniosek}
\begin{proof}
Dla $u\in C_c^\infty(Q)$ i ustalonych współrzędnych
poza $x_j$ funkcja znika w otoczeniu dolnego końca
odcinka $I_j=(a_j,a_j+L)$. Mamy więc
\[
 u(x)=\int_{a_j}^{x_j}\partial_j u(x_1,\ldots,t,\ldots,x_d)
 \,\mathrm d\lambda_1(t).
\]
Z nierówności Höldera wynika
\[
 |u(x)|^p\le L^{p-1}\int_{I_j}
 |\partial_j u(x_1,\ldots,t,\ldots,x_d)|^p
 \,\mathrm d\lambda_1(t).
\]
Dla $p=1$ jest to nierówność trójkąta dla całki.
Całkowanie względem $x_j$ daje dodatkowy czynnik $L$,
a całkowanie względem pozostałych współrzędnych daje
z Fubiniego $\|u\|_p^p\le L^p\|\partial_j u\|_p^p$
(dla $d=1$ tylko pierwsze całkowanie).
Bierzemy $p$-ty pierwiastek.

Jeżeli $u\in W^{1,p}_0(Q)$, to istnieje ciąg
$\varphi_n\in C_c^\infty(Q)$ zbieżny do $u$ w $W^{1,p}$.
W szczególności $\varphi_n\to u$ i
$\partial_j\varphi_n\to\partial_j u$ w $L^p$.
Ciągłość normy pozwala przejść do granicy
w oszacowaniu dla $\varphi_n$.

Z definicji normy mamy
$\|u\|_{W^{1,p}}^p=\|u\|_p^p+|u|_{1,p}^p$.
Poprzednia nierówność daje $\|u\|_p^p\le L^p|u|_{1,p}^p$.
Stąd otrzymujemy podane dwustronne oszacowanie.
Półnorma $|u|_{1,p}$ jest dodatnio określona na
$W^{1,p}_0(Q)$ dzięki temu oszacowaniu, a jednorodność
i nierówność trójkąta wynikają z normy skończonego
produktu przestrzeni $L^p$. Jest więc normą.
\end{proof}

\section{Reguły rachunkowe w przestrzeniach Sobolewa}
\label{sec:pde9-calculus}

\begin{twierdzenie}[Mnożenie przez funkcję gładką i iloczyn]
\label{thm:pde9-products}
Niech $\Omega\subset\R^d$ będzie otwarta, $1\le p\le\infty$
i $u\in W^{1,p}(\Omega)$.
Jeśli $\phi\in C^1(\Omega)$ oraz $\phi,\nabla\phi\in L^\infty(\Omega)$,
to $\phi u\in W^{1,p}(\Omega)$ i
$\partial_j(\phi u)=\phi\partial_j u+u\partial_j\phi$.
Ponadto jeśli $v\in W^{1,q}(\Omega)$, gdzie
$1/p+1/q=1$ z konwencją $1/\infty=0$,
to $uv\in W^{1,1}(\Omega)$ i
$\partial_j(uv)=v\partial_j u+u\partial_j v$.
W szczególności $W^{1,\infty}(\Omega)$ jest algebrą
względem mnożenia punktowego klas funkcji, z tą samą
regułą Leibniza.
\end{twierdzenie}
\begin{proof}
Ustalmy $\varphi\in C_c^\infty(\Omega)$ i wybierzmy
otwarty $U\Subset\Omega$ zawierający jej nośnik.
Dla $p=\infty$ funkcja $u$ należy lokalnie do $W^{1,1}$
przez oszacowanie na zbiorach skończonej miary.
Dla każdego rozważanego $p$ możemy zatem użyć lokalnych
wygładzeń $u_\varepsilon\to u$ w $W^{1,1}(U)$ na mocy
twierdzenia~\ref{thm:pde9-mollify} (dla $p<\infty$
korzystamy też z nierówności Höldera).
Iloczyn $\phi u_\varepsilon$ jest klasy $C^1$,
więc klasyczne całkowanie przez części daje
\[
 \int_U\phi u_\varepsilon\partial_j\varphi\,\mathrm d\lambda_d
 =-\int_U(\phi\partial_j u_\varepsilon
          +u_\varepsilon\partial_j\phi)\varphi\,\mathrm d\lambda_d.
\]
Nie podstawiamy $\phi\varphi$ bezpośrednio jako testu
dystrybucyjnego, ponieważ $\phi$ jest tylko klasy $C^1$.
Ograniczoność $\phi,\partial_j\phi$ i zbieżność
w $W^{1,1}(U)$ uzasadniają przejście do granicy
w każdej całce. Dostajemy żądaną tożsamość testową,
a globalnie
\[
 \begin{aligned}
 \|\phi u\|_p&\le\|\phi\|_\infty\|u\|_p,\\
 \|\partial_j(\phi u)\|_p
 &\le\|\phi\|_\infty\|\partial_j u\|_p+
      \|\partial_j\phi\|_\infty\|u\|_p .
 \end{aligned}
\]
Obie nierówności obowiązują również dla $p=\infty$.

Dla iloczynu $uv$ wybieramy wspólne lokalne wygładzenia
$u_\varepsilon,v_\varepsilon$ na otoczeniu $\overline U$.
Gdy $p,q<\infty$, dla $a=u$ lub $a=\partial_j u$
oraz $b=v$ lub $b=\partial_j v$ odpowiednie wygładzenia
$a_\varepsilon,b_\varepsilon$ zbiegają do $a,b$
w $L^p(U)$ i $L^q(U)$. Z nierówności Höldera mamy
\[
 \|a_\varepsilon b_\varepsilon-ab\|_{L^1(U)}
 \le\|a_\varepsilon-a\|_{L^p(U)}\|b_\varepsilon\|_{L^q(U)}
    +\|a\|_{L^p(U)}\|b_\varepsilon-b\|_{L^q(U)}\to0.
\]
Pierwszy składnik dąży do zera, ponieważ zbieżny
ciąg norm $\|b_\varepsilon\|_q$ jest ograniczony
dla małych $\varepsilon$.

Gdy $p=1,q=\infty$, to $\|b_\varepsilon\|_{L^\infty(U)}
\le\|b\|_{L^\infty(\Omega)}$ z nieujemności i normalizacji
jądra oraz wzoru na wygładzone pochodne.
Mamy $b_\varepsilon\to b$ w $L^1(U)$; dla pewnego
ciągu $\varepsilon_n\downarrow0$ wszystkie skończenie
liczne funkcje $v_{\varepsilon_n},\partial_jv_{\varepsilon_n}$
zbiegają do odpowiednich granic prawie wszędzie.
Korzystamy z faktu, że ze zbieżności w $L^1$ można
wybrać podciąg zbieżny prawie wszędzie (wybieramy
błędy norm sumowalne i stosujemy twierdzenie Tonellego
do sumy ich modułów).
Dla takiego ciągu mamy
\[
 \|a_{\varepsilon_n}b_{\varepsilon_n}-ab\|_{L^1(U)}
 \le \|b\|_\infty\|a_{\varepsilon_n}-a\|_{L^1(U)}
      +\int_U|a||b_{\varepsilon_n}-b|\,\mathrm d\lambda_d .
\]
Druga całka dąży do zera z twierdzenia o zbieżności
zmajoryzowanej, z majorantą $2\|b\|_\infty|a|$.
Przypadek $p=\infty,q=1$ otrzymujemy zamieniając
role $u$ i $v$.

Stosując te oszacowania do $uv$, $v\partial_j u$
i $u\partial_j v$, przechodzimy do granicy w tożsamości
testowej dla gładkiego iloczynu $u_{\varepsilon_n}v_{\varepsilon_n}$.
Globalnie $uv,v\partial_j u,u\partial_j v\in L^1(\Omega)$
dzięki nierówności Höldera. To dowodzi tezy.
Jeśli zaś $u,v\in W^{1,\infty}(\Omega)$, to lokalnie
możemy zastosować udowodniony przypadek $p=1,q=\infty$,
ponieważ na każdym $U\Subset\Omega$ mamy $u|_U\in W^{1,1}(U)$.
Otrzymana tożsamość dystrybucyjna obowiązuje na całej
$\Omega$, gdyż każdy test ma zwarty nośnik w pewnym takim $U$.
Wszystkie trzy iloczyny są globalnie ograniczone istotnie,
co daje przynależność do $W^{1,\infty}(\Omega)$.
\end{proof}

\begin{twierdzenie}[Reguła łańcuchowa dla funkcji klasy $C^1$ o ograniczonej pochodnej]
\label{thm:pde9-chain}
Niech $1\le p\le\infty$, $u\in W^{1,p}(\Omega)$
będzie rzeczywista, a $F\in C^1(\R)$ spełnia
$F(0)=0$ oraz $\|F'\|_\infty<\infty$.
Wtedy $F\circ u\in W^{1,p}(\Omega)$ i
\[
 \partial_j(F(u))=F'(u)\partial_j u
 \quad\text{prawie wszędzie w }\Omega.
\]
Jeżeli $\lambda_d(\Omega)<\infty$ lub $p=\infty$,
warunek $F(0)=0$ można pominąć.
\end{twierdzenie}
\begin{proof}
Niech $L=\|F'\|_\infty$. Z podstawowego twierdzenia
rachunku całkowego mamy $|F(t)-F(s)|\le L|t-s|$.
Przy $F(0)=0$ dostajemy $\|F(u)\|_p\le L\|u\|_p$.
Gdy $p<\infty$ i $\lambda_d(\Omega)<\infty$,
dodatkowa stała $F(0)$ ma normę
$|F(0)|\lambda_d(\Omega)^{1/p}$, a dla $p=\infty$
niezależnie od miary mamy $\|F(u)\|_\infty
\le |F(0)|+L\|u\|_\infty$.
W każdym z dopuszczalnych przypadków $F(u)\in L^p$.
Ponadto $\|F'(u)\partial_j u\|_p\le L\|\partial_j u\|_p$.

Dla $p<\infty$ ustalmy $U\Subset\Omega$ zawierający
nośnik rozważanego testu. Z twierdzenia
~\ref{thm:pde9-mollify} wybieramy ciąg wygładzeń
$u_n\to u$ w $W^{1,p}(U)$, a po przejściu
do podciągu także $u_n\to u$ prawie wszędzie na $U$.
Wybór podciągu uzasadnia sumowalność odpowiednio
wybranych $p$-tych potęg błędów w $L^p$ i Tonelli.
Z oszacowania Lipschitzowskiego wynika
$\|F(u_n)-F(u)\|_{L^p(U)}\le L\|u_n-u\|_{L^p(U)}\to0$.
Klasyczna reguła łańcuchowa dla $F(u_n)\in C^1(U)$
daje rozkład
\[
 \begin{aligned}
 \partial_j(F(u_n))-F'(u)\partial_j u
 &=F'(u_n)(\partial_j u_n-\partial_j u)\\
 &\quad +(F'(u_n)-F'(u))\partial_j u .
 \end{aligned}
\]
Norma pierwszego składnika jest nie większa niż
$L\|\partial_j u_n-\partial_j u\|_{L^p(U)}\to0$.
Drugi składnik dąży do zera prawie wszędzie
dzięki ciągłości $F'$; jego $p$-ta potęga modułu
jest ograniczona przez $(2L)^p|\partial_j u|^p$.
Twierdzenie o zbieżności zmajoryzowanej daje zatem
zbieżność w $L^p(U)$. Przechodząc do granicy
w całkowaniu przez części przeciw testowi,
uzyskujemy tożsamość słabej pochodnej na $\Omega$,
a wcześniejsze oszacowania zapewniają globalną
przynależność do $W^{1,p}$.

Dla $p=\infty$ stosujemy już udowodniony przypadek
$p=1$ na każdym $U\Subset\Omega$, gdzie $u$ i jej
pochodne należą do $L^1(U)$ i gdzie można pominąć
warunek $F(0)=0$. Tożsamości testowe obowiązują na całej
$\Omega$, a globalne oszacowania w $L^\infty$
dają tezę. Nie korzystamy ze zbieżności wygładzeń
w normie $W^{1,\infty}$.
\end{proof}

\begin{lemat}[Gradient znika na zbiorach poziomicowych]
\label{lem:pde9-levelsets}
Jeśli $u\in W^{1,p}(\Omega)$ jest rzeczywista,
$1\le p\le\infty$, i $c\in\R$, to
$\nabla u=0$ prawie wszędzie na $\{x\in\Omega:u(x)=c\}$.
\end{lemat}
\begin{proof}
Niech $\theta\in C^\infty(\R)$, $0\le\theta\le1$,
będzie równa $0$ na $[-1,1]$ i $1$ poza $(-2,2)$.
Taką funkcję otrzymujemy odejmując od $1$ gładką
funkcję odcinającą, skonstruowaną jak w dowodzie
twierdzenia~\ref{thm:pde8-density}.
Połóżmy $h_n(t)=\theta(n(t-c))$ i
$F_n(t)=\int_0^t h_n(s)\,\mathrm d\lambda_1(s)$,
gdzie dla $t<0$ całkę rozumiemy jako
$-\int_t^0 h_n(s)\,\mathrm d\lambda_1(s)$.
Wtedy $F_n(0)=0$, $F_n'=h_n$, $|F_n(t)|\le|t|$
oraz $|F_n(t)-t|\le\min\{|t|,4/n\}$,
ponieważ $1-h_n$ znika poza przedziałem
$(c-2/n,c+2/n)$ i ma wartości w $[0,1]$.

Dla $p<\infty$ z twierdzenia o zbieżności zmajoryzowanej
mamy $F_n(u)\to u$ w $L^p(\Omega)$,
z majorantą $|u|^p$ dla $p$-tych potęg błędów,
także na zbiorach nieskończonej miary.
Twierdzenie~\ref{thm:pde9-chain} daje
$\partial_j F_n(u)=h_n(u)\partial_j u$.
Ponieważ $h_n(t)\to\chi_{\R\setminus\{c\}}(t)$ dla
każdego $t$, otrzymujemy przez zbieżność zmajoryzowaną
\[
 h_n(u)\partial_j u\longrightarrow
 \chi_{\{u\ne c\}}\partial_j u \quad\text{w }L^p(\Omega).
\]
Rzeczywiście $p$-ta potęga modułu różnicy jest
ograniczona przez $|\partial_j u|^p$.
Z zamkniętości pojedynczej pochodnej
(wniosek~\ref{cor:pde9-closed}) dostajemy
$\partial_j u=\chi_{\{u\ne c\}}\partial_j u$.
Dla $p=\infty$ stosujemy przypadek $p=1$
na przeliczalnym pokryciu $\Omega$ relatywnie zwartymi
otwartymi podzbiorami (sekcja~\ref{sec:pde9-def}).
Przeliczalna suma zbiorów wyjątkowych ma miarę zero,
więc teza zachodzi na całej $\Omega$.
\end{proof}

\begin{wniosek}[Obcięcia, części dodatnie i moduł]
\label{cor:pde9-truncation}
Niech $u\in W^{1,p}(\Omega)$ będzie rzeczywista,
$1\le p\le\infty$.
Wtedy $u^+=\max(u,0)$, $u^-=\max(-u,0)$ oraz $|u|$
należą do $W^{1,p}(\Omega)$ i prawie wszędzie
\[
 \nabla u^+=\chi_{\{u>0\}}\nabla u,\qquad
 \nabla u^-=-\chi_{\{u<0\}}\nabla u,\qquad
 \nabla|u|=\operatorname{sgn}(u)\nabla u .
\]
Dla $M>0$ obcięcie $T_M(u)=\max(-M,\min(u,M))$
również należy do $W^{1,p}(\Omega)$ i
$\nabla T_M(u)=\chi_{\{|u|<M\}}\nabla u$ prawie wszędzie.
Jeżeli $p<\infty$, to $T_M(u)\to u$ w $W^{1,p}(\Omega)$
przy $M\to\infty$.
\end{wniosek}
\begin{proof}
Ustalmy jednowymiarowe nieujemne jądro wygładzające
$\eta\in C_c^\infty((-1,1))$ o całce $1$
i $\eta_n(t)=n\eta(nt)$.
Niech $G$ oznacza funkcję $t\mapsto t^+$
albo $t\mapsto T_M(t)$.
Obie funkcje są lipschitzowskie ze stałą $1$,
mają wartość $0$ w zerze i są absolutnie ciągłe
na każdym przedziale ograniczonym.
Ich słabe pochodne na $\R$ to odpowiednio
$g=\chi_{(0,\infty)}$ i $g=\chi_{(-M,M)}$,
co wynika z jednowymiarowej charakterystyki
~\ref{thm:pde9-ac}, stosowanej lokalnie.
Połóżmy
\[
 G_n(t)=(\eta_n*G)(t)-(\eta_n*G)(0).
\]
Z twierdzenia~\ref{thm:pde7-mollify} mamy
$G_n\in C^\infty(\R)$ i $G_n'=\eta_n*g$,
zatem $0\le G_n'\le1$ i $G_n(0)=0$.
Ponadto $|G_n(t)|\le|t|$ oraz
\[
 |(\eta_n*G)(t)-G(t)|
 \le\int_{\R}\eta_n(s)|s|\,\mathrm d\lambda_1(s)\le1/n.
\]
Stąd $|G_n(t)-G(t)|\le2/n$.
Poza skończonym zbiorem punktów załamania $G$
funkcja $g$ jest lokalnie stała, więc
$G_n'(t)=g(t)$ dla dostatecznie dużych $n$.

Dla $p<\infty$ twierdzenie~\ref{thm:pde9-chain} daje
$\partial_j G_n(u)=G_n'(u)\partial_j u$.
Zbieżność $G_n(u)\to G(u)$ w $L^p$ wynika
z majoranty $(2|u|)^p$ dla $p$-tych potęg błędów.
Poza przeciwobrazami punktów załamania mamy
$G_n'(u)\to g(u)$, zaś na tych przeciwobrazach
$\partial_j u=0$ prawie wszędzie z lematu
~\ref{lem:pde9-levelsets}.
Zatem $G_n'(u)\partial_j u\to g(u)\partial_j u$
prawie wszędzie na całej $\Omega$.
Ponieważ $0\le G_n',g\le1$, $p$-ta potęga modułu
różnicy jest ograniczona przez $|\partial_j u|^p$.
Zbieżność zmajoryzowana i wniosek~\ref{cor:pde9-closed}
dają $\partial_j G(u)=g(u)\partial_j u$.
Wszystkie pochodne są w $L^p$ i $|G(u)|\le|u|$.
Dla $p=\infty$ stosujemy te same tożsamości lokalnie
z $p=1$, a globalnie korzystamy z oszacowań
$|G(u)|\le|u|$ i $|\partial_j G(u)|\le|\partial_j u|$.
Wzór dla $u^-$ otrzymujemy stosując wynik dla części
dodatniej do $-u$, a dla $|u|=u^++u^-$ przez liniowość.
Przyjmujemy $\operatorname{sgn}(0)=0$.

Wreszcie dla $p<\infty$ i $M\to\infty$
mamy $|T_M(u)-u|\le|u|\chi_{\{|u|>M\}}$
oraz
$\partial_j T_M(u)-\partial_j u
=-\chi_{\{|u|\ge M\}}\partial_j u$ prawie wszędzie.
Oba błędy dążą do zera w $L^p$ dzięki zbieżności
zmajoryzowanej, odpowiednio przez $|u|^p$
i $|\partial_j u|^p$.
Skończona suma norm pochodnych i funkcji daje tezę.
\end{proof}

\begin{uwaga}[Zakres reguły łańcuchowej]
Twierdzenie~\ref{thm:pde9-chain} dotyczy funkcji
klasy $C^1$ o ograniczonej pochodnej, a nie bezpośrednio
wszystkich funkcji lipschitzowskich. Powyższy argument
dla obcięć wykorzystuje w istotny sposób skończoną
liczbę punktów załamania i lemat o zbiorach poziomicowych.
Nie należy bez dodatkowego uzasadnienia zastępować
$F'(u)$ dowolnym reprezentantem pochodnej ogólnej
funkcji lipschitzowskiej.
\end{uwaga}

\section{Przedłużanie z półprzestrzeni i rola brzegu}
\label{sec:pde9-extension}

Najpierw uzasadnimy aproksymację niezależnie od przedłużania.
Niech $P=\R^{d-1}\times(0,\infty)$, $1\le p<\infty$,
$u\in W^{k,p}(P)$, $k\in\N_0$. Dla $d=1$ pomijamy $x'$.
Dla $h>0$ określamy $u_h(x',t)=u(x',t+h)$ na
$P_h=\R^{d-1}\times(-h,\infty)$. Zmiana zmiennych w
tożsamości testowej daje $\partial^\alpha u_h(x',t)
=(\partial^\alpha u)(x',t+h)$ dla $|\alpha|\le k$.
Niech $f_\alpha$ będzie przedłużeniem przez zero
$\partial^\alpha u$ jako funkcji z $L^p(P)$.
Ciągłość przesunięć w $L^p$ daje
\[
 \|\partial^\alpha(u_h-u)\|_{L^p(P)}
 \le \|f_\alpha(\cdot+he_d)-f_\alpha\|_{L^p(\R^d)}
 \longrightarrow0.
\]
Nie zakładamy, że przedłużenia $f_\alpha$ są słabymi
pochodnymi $f_0$. Analogicznie niech $g_{h,\alpha}$
będą oddzielnymi przedłużeniami przez zero pochodnych
$u_h$ z $P_h$ jako funkcje $L^p$.
Dla $0<\varepsilon<h/2$ funkcja
$w_{h,\varepsilon}=\rho_\varepsilon*g_{h,0}$
jest gładka na $\R^d$, a lokalna reguła splotu
(twierdzenie~\ref{thm:pde9-mollify}) daje na $P$
$\partial^\alpha w_{h,\varepsilon}
=\rho_\varepsilon*g_{h,\alpha}$, ponieważ nośnik jądra
przesuniętego do punktu $P$ mieści się w $P_h$.
Z twierdzenia o przybliżonej jedynce
(rozdział~\ref{chap:pde-distributions}) mamy
\[
 \|\partial^\alpha(w_{h,\varepsilon}-u_h)\|_{L^p(P)}
 \le \|\rho_\varepsilon*g_{h,\alpha}-g_{h,\alpha}\|_{L^p(\R^d)}
 \longrightarrow0
\]
dla ustalonego $h$. Ponieważ wielowskaźników jest
skończenie wiele, wybieramy $h_n\downarrow0$ oraz
$0<\varepsilon_n<h_n/2$ z błędem w $W^{k,p}(P)$
między $w_{h_n,\varepsilon_n}$ a $u_{h_n}$ mniejszym
od $1/n$. Otrzymujemy globalną zbieżność gładkich
restrykcji do $u$ w $W^{k,p}(P)$.

\begin{twierdzenie}[Odbicie Sobolewa, rząd pierwszy]
\label{thm:pde9-reflection}
Niech $\R^d_+=\{(x',t)\in\R^{d-1}\times\R:t>0\}$,
$1\le p<\infty$. Istnieje liniowy, ograniczony operator
$E:W^{1,p}(\R^d_+)\to W^{1,p}(\R^d)$ taki, że
$(Eu)|_{\R^d_+}=u$. Można przyjąć
\[
 (Eu)(x',t)=\begin{cases}u(x',t),&t>0,\\
 u(x',-t),&t<0.\end{cases}
\]
Wówczas $\|Eu\|_{W^{1,p}(\R^d)}=2^{1/p}
\|u\|_{W^{1,p}(\R^d_+)}$.
\end{twierdzenie}
\begin{proof}
Wartości na $t=0$ są nieistotne dla klas $L^p$,
a odbicie zachowuje zbiory miary zero, więc wzór jest
dobrze określony i liniowy.
Dla restrykcji funkcji globalnie gładkiej zwykłe
różniczkowanie na obu półprzestrzeniach daje
$\partial_t Eu(x',t)=\operatorname{sgn}(t)
(\partial_tu)(x',|t|)$; pochodne styczne odbijają się
bez zmiany znaku. Całkujemy przez części przeciw
testowi z $C_c^\infty(\R^d)$ osobno na obu
półprzestrzeniach, ograniczając się do prostopadłościanu,
w którego wnętrzu mieści się nośnik testu.
Dla pochodnej normalnej składniki brzegowe na $t=0$
to całki z $-u(x',0)\varphi(x',0)$ i
$+u(x',0)\varphi(x',0)$ (dla $d=1$ same wartości).
Znaki pochodzą od przeciwnych normalnych zewnętrznych.
Składniki znoszą się; dla pochodnych stycznych nie
występują. Otrzymujemy tożsamości słabych pochodnych.
Podstawienie $s=-t$ daje dla funkcji i każdej pochodnej
podwojenie $p$-tej potęgi normy. Po zsumowaniu
$\|Eu\|_{W^{1,p}}=2^{1/p}\|u\|_{W^{1,p}}$.

Dla ogólnego $u$ wybieramy gładkie restrykcje $w_n$
z poprzedniej aproksymacji. Równość norm dla różnic
daje ciąg Cauchy'ego $Ew_n$ w $W^{1,p}(\R^d)$,
więc z zupełności istnieje jego granica $v$.
Jednocześnie sam wzór odbicia określa operator na $L^p$
o normie $2^{1/p}$, więc $Ew_n\to Eu$ w $L^p$.
Jednoznaczność granicy daje $v=Eu$ prawie wszędzie;
równość norm przechodzi do granicy.
\end{proof}

\begin{stwierdzenie}[Odbicia wyższego rzędu na półprzestrzeni]
\label{prop:pde9-highreflection}
Niech $k\in\N$, $1\le p<\infty$. Istnieje ograniczony,
liniowy operator $E_k:W^{k,p}(\R^d_+)\to W^{k,p}(\R^d)$,
który jest prawym odwrotnym do ograniczenia na $\R^d_+$.
\end{stwierdzenie}
\begin{proof}
Wybieramy liczby $a_1,\ldots,a_k$ spełniające
$\sum_{j=1}^k a_j(-j)^m=1$ dla $m=0,\ldots,k-1$.
Macierz układu jest macierzą Vandermonde'a o różnych
węzłach $-1,\ldots,-k$. Zależność liniowa jej wierszy
dawałaby wielomian stopnia co najwyżej $k-1$
z $k$ różnymi pierwiastkami. Układ ma więc dokładnie
jedno rozwiązanie. Definiujemy $E_ku=u$ dla $t>0$,
a dla $t<0$
\[
 (E_ku)(x',t)=\sum_{j=1}^k a_j u(x',-jt).
\]
Zmiany zmiennych są odwracalne i zachowują zbiory
miary zero, więc wzór jest liniowy na klasach $L^p$.
Dla gładkiej restrykcji $u$ pochodne klasyczne po
dolnej stronie mają postać
\[
 \partial_t^m\partial_{x'}^\beta E_ku(x',t)
 =\sum_{j=1}^k a_j(-j)^m
 (\partial_t^m\partial_{x'}^\beta u)(x',-jt).
\]
Jeżeli $m+|\beta|\le k-1$, to wartość graniczna
przy $t\uparrow0$ jest równa odpowiedniej pochodnej
od góry dzięki układowi na $a_j$.
Indukcyjnie dla $r=0,\ldots,k-1$ całkujemy przez części
każdą zidentyfikowaną pochodną rzędu $r$ jak w
dowodzie odbicia pierwszego rzędu. Dopasowanie znosi
składniki brzegowe i identyfikuje słabą pochodną
rzędu $r+1$ z klasyczną pochodną na obu częściach.
Nie wymagamy dopasowania pochodnych rzędu $k$.

Niech $C_m=\sum_{j=1}^k |a_j|j^{m-1/p}$.
Nierówność trójkąta i podstawienie $s=-jt$ dają
\[
 \|\partial_t^m\partial_{x'}^\beta E_ku\|_{L^p(t<0)}
 \le C_m\|\partial_t^m\partial_{x'}^\beta u\|_{L^p(P)}.
\]
Istotnie, w $p$-tej potędze normy pojedynczego składnika
różniczkowanie daje czynnik $j^{mp}$, a jakobian $j^{-1}$.
Dodając normy po górnej stronie i sumując, otrzymujemy
\[
 \|E_ku\|_{W^{k,p}(\R^d)}
 \le \max_{0\le m\le k}(1+C_m^p)^{1/p}
 \|u\|_{W^{k,p}(P)}.
\]
Dla ogólnego $u$ stosujemy poprzednią niezależną
aproksymację rzędu $k$. Oszacowanie różnic i zupełność
dają granicę przedłużeń w $W^{k,p}(\R^d)$.
Oszacowanie dla $m=0,\beta=0$ działa na samych
funkcjach $L^p$, więc granica jest równa funkcji
określonej przez jawny wzór na $E_k$. Ograniczenie do $P$
jest równe $u$, a oszacowanie przechodzi do granicy.
\end{proof}

\begin{uwaga}[Obszary Lipschitzowskie]
Jeśli ograniczony obszar $\Omega$ ma brzeg lokalnie
będący wykresem funkcji Lipschitzowskiej, można prostować
brzeg dwulipschitzowskimi zmianami współrzędnych,
zastosować odbicie pierwszego rzędu i skleić lokalne
przedłużenia podziałem jedności. Otrzymujemy ograniczony
operator $W^{1,p}(\Omega)\to W^{1,p}(\R^d)$.
Przypadek $W^{k,p}$ dla $k>1$ na dowolnym obszarze Lipschitzowskim wymaga bardziej subtelnej konstrukcji:
proste prostowanie Lipschitzowskie nie zachowuje ogólnie
drugich pochodnych w $L^p$. Pełną teorię przedłużania
i śladów rozwiniemy w następnym rozdziale.
\end{uwaga}

\begin{lemat}[Gęstość funkcji gładkich aż do płaskiego brzegu]
\label{lem:pde9-density-halfspace-slab}
Niech $1\le p<\infty$ oraz $d\ge1$. Dla
$G=\R^{d-1}\times(0,\infty)$ albo
$G=\R^{d-1}\times(0,1)$ każdą funkcję $u\in W^{1,p}(G)$
można przybliżyć w normie $W^{1,p}(G)$ restrykcjami funkcji
z $C^\infty(\R^d)$. Można przy tym wybrać przybliżenia
z $C_c^\infty(\R^d)$.
\end{lemat}
\begin{proof}
Dla półprzestrzeni korzystamy z twierdzenia
~\ref{thm:pde9-reflection}. Dla warstwy wybieramy
gładką funkcję $\theta(t)$ równą $1$ dla $t\le1/3$
i $0$ dla $t\ge2/3$, o wartościach w $[0,1]$.
Rozkładamy $u=u_0+u_1$, gdzie $u_0=\theta u$,
$u_1=(1-\theta)u$. Reguła iloczynu
(twierdzenie~\ref{thm:pde9-products}) daje ograniczone
liniowe odwzorowania $u\mapsto u_i$ w $W^{1,p}(G)$.

Przedłużamy $u_0$ przez zero dla $t\ge1$ do $P$.
Dla testu na $P$ mnożymy go przez gładką funkcję
zmiennej $t$, równą $1$ dla $t\le2/3$ i $0$ dla
$t\ge5/6$. Otrzymujemy test na $G$. Tam, gdzie
zmienia się odcięcie, $u_0$ i jego słabe pochodne są
zerowe dzięki lokalności pochodnych. Tożsamość testowa
na $G$ staje się zatem tożsamością na $P$ dla
przedłużeń $u_0$ i jego pochodnych przez zero.
Normy pozostają równe; nie zakładamy zerowego śladu
przy $t=0$. Dla $u_1$ stosujemy zmianę $s=1-t$
i tę samą konstrukcję względem $s$. Zmiana zachowuje
normę, a pochodna normalna zmienia znak, co wynika
z podstawienia w tożsamości testowej.
Obie uzyskane funkcje z $W^{1,p}(P)$ przedłużamy
przez odbicie. Drugie przedłużenie przenosimy
z powrotem do zmiennej $t=1-s$ i dodajemy do pierwszego.
Na $G$ suma jest równa $u$. Nierówność trójkąta,
oszacowania iloczynów i równość norm odbicia dają
ograniczony liniowy operator przedłużania.
Nie korzystamy z nieudowodnionej jeszcze własności
absolutnej ciągłości na prawie wszystkich prostych.

W obu przypadkach niech $v=Eu$. Gęstość
$C_c^\infty(\R^d)$ w $W^{1,p}(\R^d)$ udowodniona
w sekcji~\ref{sec:pde9-density} daje $v_n\to v$
w tej normie. Ograniczenie nie zwiększa normy,
więc $\|v_n|_G-u\|_{W^{1,p}(G)}
\le\|v_n-v\|_{W^{1,p}(\R^d)}\to0$.
\end{proof}


\begin{uwaga}[Aproksymacja gładka nie oznacza zerowego śladu]
\label{rem:pde-v42-density}
Twierdzenie Meyersa--Serrina~\ref{thm:pde9-MeyersSerrin}
nie orzeka, że $C_c^\infty(\Omega)$ jest gęste w całej
przestrzeni $W^{1,p}(\Omega)$. Na przykład funkcja
$u\equiv1$ na ograniczonym obszarze Lipschitzowskim
należy do $W^{1,p}(\Omega)$, ale jej ślad jest równy $1$,
a więc nie należy do $W_0^{1,p}(\Omega)$ na mocy
twierdzenia~\ref{thm:pde10-trace-zero}.
Są to trzy różne twierdzenia: aproksymacja wewnętrzna,
gęstość gładkich restrykcji aż do brzegu oraz
charakterystyka jądra śladu. Ten ostatni rezultat
zostanie udowodniony w rozdziale~\ref{chap:pde-embeddings}.
\end{uwaga}

\section{Ślady na płaskim brzegu: pierwszy wynik}
\label{sec:pde9-traces}

\begin{twierdzenie}[Oszacowanie śladu na warstwie]
\label{thm:pde9-trace-slab}
Niech $S=\R^{d-1}\times(0,1)$ i $1\le p<\infty$.
Istnieje dokładnie jeden ograniczony operator liniowy
$\operatorname{Tr}_0:W^{1,p}(S)\to L^p(\R^{d-1})$
zgodny z ograniczeniem do $t=0$ dla restrykcji funkcji
gładkich na otoczeniu $\overline S$, należących do $W^{1,p}(S)$.
Dla $d=1$ przestrzeń wartości jest równa $\R$ z normą
modułu (ślad w punkcie). Zachodzi
\[
 \|\operatorname{Tr}_0u\|_{L^p(\R^{d-1})}
 \le \|u\|_{L^p(S)}+\|\partial_tu\|_{L^p(S)}.
\]
Analogiczne twierdzenie zachodzi dla śladu $\operatorname{Tr}_1$
na $t=1$.
\end{twierdzenie}
\begin{proof}
Dla funkcji $h$ klasy $C^1$ na $[0,1]$ i $s\in(0,1)$
z podstawowego twierdzenia rachunku całkowego otrzymujemy
\[
 |h(0)|\le |h(s)|+\int_0^1|h'(t)|\,\mathrm d\lambda_1(t).
\]
Całkujemy po $s$ i stosujemy nierówność Höldera
(na przedziale miary $1$; dla $p=1$ nie jest potrzebna):
\[
 |h(0)|\le\|h\|_{L^p(0,1)}+\|h'\|_{L^p(0,1)}.
\]
Dla gładkiej restrykcji $u$ stosujemy to do
$h(t)=u(x',t)$. Z twierdzenia Tonellego funkcje
\[
 A(x')=\left(\int_0^1 |u(x',t)|^p\,\mathrm d\lambda_1(t)\right)^{1/p},
 \qquad
 B(x')=\left(\int_0^1 |\partial_tu(x',t)|^p\,\mathrm d\lambda_1(t)\right)^{1/p}
\]
są mierzalne, skończone prawie wszędzie, a ich normy
$L^p(\R^{d-1})$ są równe odpowiednim normom na $S$.
Nierówność punktowa $|u(x',0)|\le A(x')+B(x')$
oraz nierówność trójkąta w $L^p$ dają żądane
oszacowanie. Dla $d=1$ wystarcza bezpośrednio
oszacowanie jednowymiarowe.

Z lematu~\ref{lem:pde9-density-halfspace-slab} dla dowolnego
$u\in W^{1,p}(S)$ istnieje ciąg gładkich restrykcji $u_n$
zbieżny do $u$ w tej przestrzeni. Oszacowanie zastosowane
do $u_n-u_m$ pokazuje, że $u_n(\cdot,0)$ jest ciągiem
Cauchy'ego w przestrzeni Banacha $L^p(\R^{d-1})$
(dla $d=1$ w $\R$). Definiujemy $\operatorname{Tr}_0u$
jako jego granicę. Jeżeli $v_n$ jest drugim ciągiem
aproksymującym, to $u_n-v_n\to0$ w $W^{1,p}$, a więc
różnica ich śladów dąży do zera w $L^p$. Definicja
nie zależy od wyboru ciągu. Aproksymując kombinacje
liniowe przez te same kombinacje aproksymantów,
otrzymujemy liniowość. Przejście do granicy w
oszacowaniu dla $u_n$ daje oszacowanie dla $u$.
Dla gładkiej restrykcji można użyć ciągu stałego,
więc operator zgadza się z ograniczeniem klasycznym.
Każdy inny ograniczony operator o tej własności
musi dawać na $u$ tę samą granicę, co dowodzi
jednoznaczności. Zmiana zmiennej $t\mapsto1-t$
zachowuje normy funkcji i modułów pochodnych normalnych,
co daje twierdzenie dla $\operatorname{Tr}_1$.
\end{proof}

Jeśli $u\in W^{1,p}_0(S)$, oba ślady są zerowe:
aproksymanty z $C_c^\infty(S)$ znikają w otoczeniu
obu płaszczyzn brzegowych, a operatory śladu są
ciągłe. Nie dowodzimy tu jeszcze odwrotnej implikacji
dla ogólnych obszarów; por. rozdział~\ref{chap:pde-embeddings}.

\section{Perspektywa wariacyjna: pierwsze słabe rozwiązanie}
\label{sec:pde9-weakPDE}

\begin{przyklad}[Równanie Poissona na przedziale]
\label{ex:pde9-dirichlet}
Niech $I=(0,1)$, $f\in L^2(I)$ i rozważmy
$-u''=f$ na $I$ z warunkami $u(0)=u(1)=0$.
Jeśli $u$ jest gładkie i spełnia równanie, mnożymy je
przez $\varphi\in C_c^\infty(I)$ i całkujemy przez części:
\[
 \int_0^1u'(x)\varphi'(x)\,\mathrm d\lambda_1(x)
 =\int_0^1 f(x)\varphi(x)\,\mathrm d\lambda_1(x).
\]
Obie strony mają sens dla $u,\varphi\in H^1(I)$
dzięki nierówności Cauchy'ego--Schwarza.
Szukamy $u\in H_0^1(I)$ spełniającego tę tożsamość
dla każdego $\varphi\in H_0^1(I)$.
Jest to \emph{słabe zagadnienie Dirichleta}; zerowe
wartości brzegowe reprezentanta $u$ wynikają z
ciągłości śladów w wymiarze jeden
(wniosek~\ref{cor:pde9-trace-oneD}) i definicji $H_0^1$.

Wykażemy istnienie jawnie. Ponieważ $f\in L^1(I)$,
funkcja $F(x)=\int_0^x f(y)\,\mathrm d\lambda_1(y)$
jest absolutnie ciągła na $[0,1]$ i $F'=f$ prawie wszędzie.
Połóżmy
\[
 c=\int_0^1 F(s)\,\mathrm d\lambda_1(s),\qquad
 u(x)=cx-\int_0^x F(s)\,\mathrm d\lambda_1(s).
\]
Funkcja $F$ jest ciągła i ograniczona na $[0,1]$,
zatem $u$ jest klasy $C^1$ i $u'=c-F$ jest
absolutnie ciągła. Mamy $u''=-f$ prawie wszędzie,
a z jednowymiarowej charakterystyki słabych pochodnych
(twierdzenie~\ref{thm:pde9-ac}) wynika $u\in H^2(I)$.
Ponadto $u(0)=0$ i $u(1)=c-c=0$.

Uzasadnimy jeszcze $u\in H_0^1(I)$ bez zakładania
ogólnego twierdzenia o jądrze śladu. Dla
$0<\varepsilon<1/4$ wybierzmy gładkie odcięcie
$\eta_\varepsilon$ o wartościach w $[0,1]$, równe
$1$ na $[2\varepsilon,1-2\varepsilon]$, równe $0$
na $(0,\varepsilon)\cup(1-\varepsilon,1)$,
z $|\eta_\varepsilon'|\le C/\varepsilon$,
gdzie $C$ nie zależy od $\varepsilon$.
Oznaczmy $M=\sup_{[0,1]}|u'|$. Z zerowych wartości
brzegowych i całkowania pochodnej wynika
$|u(x)|\le M\min(x,1-x)$. Na zbiorze, gdzie
$\eta_\varepsilon'\ne0$, mamy więc
$|\eta_\varepsilon' u|\le 2CM$, a miara tego zbioru
jest nie większa niż $4\varepsilon$. Stąd
$\|\eta_\varepsilon' u\|_2^2\le16 C^2M^2\varepsilon$.
Pozostałe składniki błędu,
$(1-\eta_\varepsilon)u$ i $(1-\eta_\varepsilon)u'$,
dążą do zera w $L^2(I)$ na mocy twierdzenia o
zbieżności zmajoryzowanej. Reguła Leibniza daje więc
$\eta_\varepsilon u\to u$ w $H^1(I)$.
Każda funkcja $\eta_\varepsilon u$ ma zwarty nośnik
w $I$ i należy do $H_0^1(I)$ dzięki wygładzaniu
z zachowaniem zwartego nośnika
(lemat~\ref{lem:pde9-compact-product}, zastosowany do
$u$ i $\eta_\varepsilon$). Domkniętość $H_0^1(I)$
daje $u\in H_0^1(I)$.

Dla testów z $C_c^\infty(I)$ tożsamość słaba wynika
z równania $-u''=f$ w sensie dystrybucji. Dla
$\varphi\in H_0^1(I)$ wybieramy testy $\varphi_n\to\varphi$
w $H^1(I)$. Obie strony tożsamości zbiegają do
odpowiednich granic, ponieważ
\[
 \left|\int_0^1u'(\varphi_n'-\varphi')\,\mathrm d\lambda_1\right|
 \le \|u'\|_2\|\varphi_n'-\varphi'\|_2,\qquad
 \left|\int_0^1f(\varphi_n-\varphi)\,\mathrm d\lambda_1\right|
 \le \|f\|_2\|\varphi_n-\varphi\|_2.
\]
Otrzymaliśmy zatem rozwiązanie słabe.

Twierdzenie Fubiniego (stosowalne, gdyż $f\in L^1(I)$
i wszystkie użyte wagi są ograniczone) daje
$c=\int_0^1(1-y)f(y)\,\mathrm d\lambda_1(y)$ oraz
$\int_0^xF(s)\,\mathrm d\lambda_1(s)
=\int_0^x(x-y)f(y)\,\mathrm d\lambda_1(y)$.
Po rozdzieleniu pierwszej całki na $(0,x)$ i $(x,1)$
otrzymujemy równoważny wzór
\[
 u(x)=\int_0^1G(x,y)f(y)\,\mathrm d\lambda_1(y),\qquad
 G(x,y)=\begin{cases}y(1-x),&y\le x,\\x(1-y),&x<y.\end{cases}
\]
Dla $y\in I$ funkcja $G(\cdot,y)$ jest ciągła,
zeruje się na końcach i ma pochodną $1-y$ dla $x<y$
oraz $-y$ dla $x>y$. Całkując przeciw pochodnej
testu, otrzymujemy
\[
 \int_0^1\partial_xG(x,y)\varphi'(x)\,\mathrm d\lambda_1(x)
 =(1-y)\varphi(y)+y\varphi(y)=\varphi(y).
\]
Zatem $-\partial_x^2G(\cdot,y)=\delta_y$.
Dla $f\equiv1$ powyższa konstrukcja daje $F(x)=x$,
$c=1/2$ i $u(x)=x(1-x)/2$.

Jeśli $u_1,u_2$ są dwoma słabymi rozwiązaniami,
to $w=u_1-u_2\in H_0^1(I)$ spełnia jednorodną
tożsamość. Podstawiając $\varphi=w$ otrzymujemy
$\|w'\|_2^2=0$. Nierówność Poincarégo
(wniosek~\ref{cor:pde9-Poincare-zero}, $L=1$)
daje $\|w\|_2=0$, a więc jednoznaczność.
Podstawiając $\varphi=u$ dla dowolnego $f$, mamy
\[
 \|u'\|_2^2=\int_0^1fu\,\mathrm d\lambda_1
 \le\|f\|_2\|u\|_2\le\|f\|_2\|u'\|_2 .
\]
Jeśli $\|u'\|_2=0$, oszacowanie jest natychmiastowe;
w przeciwnym razie dzielimy przez tę normę.
Dostajemy $\|u'\|_2\le\|f\|_2$ i
$\|u\|_{H^1(I)}\le\sqrt{2}\|f\|_2$.
\end{przyklad}

\section{Skalowanie, krytyczne wykładniki i zadania na przyszłość}
\label{sec:pde9-scaling}

Niech $\phi\in C_c^\infty(\R^d)$, $\phi\not\equiv0$,
$a\in\R$, a dla $\varepsilon>0$ połóżmy
$u_\varepsilon(x)=\varepsilon^{-a}\phi(x/\varepsilon)$.
Dla $1\le p,q<\infty$ podstawienie $x=\varepsilon y$
o jakobianie $\varepsilon^d$ daje
\[
 \|u_\varepsilon\|_{L^q}^q
 =\varepsilon^{d-aq}\int_{\R^d}|\phi(y)|^q\,\mathrm d\lambda_d(y),
 \qquad
 \nabla u_\varepsilon(x)=\varepsilon^{-a-1}\nabla\phi(x/\varepsilon).
\]
Stąd, przy normie euklidesowej gradientu,
\[
 \|u_\varepsilon\|_{L^q}
 =\varepsilon^{d/q-a}\|\phi\|_{L^q},\qquad
 \|\nabla u_\varepsilon\|_{L^p}
 =\varepsilon^{d/p-a-1}\|\nabla\phi\|_{L^p}.
\]
Te same wykładniki obowiązują dla $q=\infty$ lub
$p=\infty$ z konwencją $1/\infty=0$, ponieważ
skalowanie argumentu zachowuje zbiory miary zero,
a mnożenie przez stałą skaluje istotne supremum.

Jeśli nierówność jednorodna
$\|u\|_{L^q}\le C\|\nabla u\|_{L^p}$ na $\R^d$
ma obowiązywać z jedną stałą $C$ dla wszystkich
$u\in C_c^\infty(\R^d)$, to po podstawieniu
$u_\varepsilon$ i podzieleniu przez odpowiednią
potęgę $\varepsilon$ dostajemy
\[
 \varepsilon^{d/q-d/p+1}\|\phi\|_{L^q}
 \le C\|\nabla\phi\|_{L^p}\qquad(\varepsilon>0).
\]
Norma $\|\phi\|_{L^q}$ jest dodatnia. Jeżeli wykładnik
po lewej jest ujemny, otrzymujemy sprzeczność przy
$\varepsilon\downarrow0$, a jeżeli dodatni,
przy $\varepsilon\to\infty$. Koniecznie więc
$d/q=d/p-1$. Dla $1\le p<d$ otrzymujemy wykładnik
\[
 p^*=\frac{dp}{d-p}.
\]
Warunek skalowania jest konieczny, lecz sam nie dowodzi
nierówności Sobolewa. W szczególności dla $p=d$ wskazuje
tylko potencjalny wykładnik $q=\infty$, nie gwarantuje
prawdziwości oszacowania w tym przypadku.

Skalowanie daje również kontrprzykład do ciągłego
zanurzenia w zbyt wysokim wykładniku nawet na dowolnym
niepustym zbiorze otwartym $\Omega\subset\R^d$.
Niech $1\le p<d$, $p^*<q\le\infty$,
$B_r(x_0)\subset\Omega$ i wybierzmy niezerową
$\phi\in C_c^\infty(B_1(0))$. Dla $0<\varepsilon<r$
połóżmy
$w_\varepsilon(x)=\varepsilon^{1-d/p}
\phi((x-x_0)/\varepsilon)$.
Wtedy $w_\varepsilon\in C_c^\infty(\Omega)$, a
\[
 \|w_\varepsilon\|_{L^p(\Omega)}=\varepsilon\|\phi\|_p,
 \qquad
 \|\partial_j w_\varepsilon\|_{L^p(\Omega)}
 =\|\partial_j\phi\|_p,\qquad
 \|w_\varepsilon\|_{L^q(\Omega)}
 =\varepsilon^{d/q-d/p+1}\|\phi\|_q.
\]
Pierwsze dwa wzory pokazują ograniczoność w
$W^{1,p}(\Omega)$ przy $\varepsilon\downarrow0$,
natomiast ostatnia norma dąży do nieskończoności,
ponieważ $d/q<d/p-1$. Wyklucza to ciągłe
zanurzenie $W^{1,p}(\Omega)\hookrightarrow L^q(\Omega)$
dla $q>p^*$, również w przypadku $q=\infty$.
Nie twierdzimy na podstawie samego tego argumentu,
że istnieje pojedynczy element $W^{1,p}(\Omega)$
nienależący do $L^q(\Omega)$. W następnym rozdziale
udowodnimy dodatnie twierdzenia o zanurzeniach
i zbadamy ich zwartość.

\section{Podsumowanie i zależności logiczne}
Wprowadziliśmy przestrzenie $W^{k,p}$ na dowolnych
niepustych zbiorach otwartych, wykazaliśmy zupełność
i zamkniętość pochodnych. Jednowymiarowa charakterystyka
przez absolutną ciągłość dała ślady w punktach.
Dla $p<\infty$ udowodniliśmy lokalne wygładzanie
i twierdzenie Meyersa--Serrina oraz ograniczone
przedłużenia z półprzestrzeni; przypadki $p=\infty$
rozpatrywaliśmy osobno tam, gdzie były dopuszczalne.
Oddzieliliśmy $W^{k,p}(\Omega)$ od domknięcia
$W^{k,p}_0(\Omega)$, którego nie utożsamiamy
na dowolnym zbiorze otwartym z jądrem operatora śladu.
Udowodniliśmy nierówności Poincarégo na kostkach,
reguły rachunkowe i oszacowanie śladu na warstwie,
a także skonstruowaliśmy pierwsze słabe rozwiązanie
równania Poissona. W rozdziale~\ref{chap:pde-embeddings}
wykorzystamy te narzędzia do ogólnych twierdzeń o
zanurzeniach, zwartości Rellicha--Kondrachova
i teorii śladów na brzegach Lipschitzowskich.


% Rozdział 10. Samodzielna część wykładu; włączana przed dodatkiem z kodem.
